% =====================================================================
% Phase 3 — L'agent PyTutor (LangChain 1.x)
% Concepts essentiels des agents ReAct avec tools et mémoire.
% Compilation : pdflatex -shell-escape phase3.tex (deux fois pour la TOC)
% Pré-requis  : pygmentize (pip install Pygments) + paquet TeX "minted"
% =====================================================================
\documentclass[11pt,a4paper]{article}

% ---------- Encodage & langue ----------
\usepackage[utf8]{inputenc}
\usepackage[T1]{fontenc}
\usepackage[french]{babel}
\usepackage{lmodern}
\usepackage{microtype}

% Glyphes Unicode utilisés dans le code Python copié-collé
\DeclareUnicodeCharacter{2713}{\textbf{V}}    % ✓
\DeclareUnicodeCharacter{2717}{\textbf{X}}    % ✗
\DeclareUnicodeCharacter{2192}{$\rightarrow$} % →
\DeclareUnicodeCharacter{2500}{-}             % ─ box drawing
\DeclareUnicodeCharacter{276F}{>}             % ❯

% ---------- Mise en page ----------
\usepackage[a4paper,margin=2.2cm]{geometry}
\usepackage{parskip}
\usepackage{xcolor}
\usepackage{enumitem}
\usepackage{booktabs}
\usepackage{hyperref}
\hypersetup{
  colorlinks=true,
  linkcolor=blue!60!black,
  urlcolor=blue!60!black,
  citecolor=blue!60!black,
  bookmarksopen=true,
  bookmarksopenlevel=2,
}

% ---------- Boîtes & callouts ----------
\usepackage[most]{tcolorbox}
\tcbset{
  enhanced,
  colback=blue!3,
  colframe=blue!50!black,
  fonttitle=\bfseries,
  arc=2pt, boxrule=0.6pt,
}

% ---------- Code avec minted ----------
\usepackage{minted}
\setminted{
  fontsize=\small,
  linenos=true,
  numbersep=6pt,
  frame=lines,
  framesep=2mm,
  bgcolor=gray!4,
  breaklines=true,
  breakanywhere=true,
  python3=true,
  tabsize=4,
}
\newminted{python}{}

% ---------- TikZ ----------
\usepackage{tikz}
\usetikzlibrary{
  arrows.meta, positioning, shapes.geometric, shapes.misc,
  fit, backgrounds, calc, trees,
  decorations.pathreplacing, decorations.markings,
}

\tikzset{
  node distance=8mm and 12mm,
  block/.style={
    rectangle, rounded corners=2pt, draw=blue!60!black, fill=blue!8,
    minimum width=42mm, minimum height=9mm, align=center, font=\small
  },
  data/.style={
    rectangle, rounded corners=1pt, draw=gray!60, fill=gray!10,
    minimum width=42mm, minimum height=8mm, align=center, font=\footnotesize\ttfamily
  },
  llm/.style={
    rectangle, rounded corners=2pt, draw=red!60!black, fill=red!8,
    minimum width=42mm, minimum height=9mm, align=center, font=\small
  },
  tool/.style={
    rectangle, rounded corners=2pt, draw=teal!60!black, fill=teal!10,
    minimum width=40mm, minimum height=9mm, align=center, font=\small
  },
  store/.style={
    cylinder, shape border rotate=90, aspect=0.25,
    draw=violet!60!black, fill=violet!10, minimum width=30mm, minimum height=12mm,
    align=center, font=\small
  },
  schema/.style={
    rectangle, rounded corners=2pt, draw=teal!60!black, fill=teal!8,
    minimum width=48mm, minimum height=9mm, align=center, font=\small
  },
  msgblock/.style={
    rectangle, rounded corners=1pt, draw=purple!60!black, fill=purple!6,
    minimum width=58mm, minimum height=7mm, align=left, font=\scriptsize\ttfamily
  },
  flow/.style={-{Stealth[length=2.2mm]}, thick, draw=blue!60!black},
  flowdata/.style={-{Stealth[length=2.2mm]}, thick, draw=gray!60, dashed},
  label/.style={font=\scriptsize\itshape, fill=white, inner sep=1pt},
}

% ---------- Titre ----------
\title{\textbf{Phase 3 — L'agent PyTutor}\\
\large Quatre concepts essentiels + cinq exercices pour aller plus loin}
\author{Cours PyTutor2}
\date{}

\begin{document}
\maketitle

\begin{abstract}
\noindent
Cette phase 3 construit un \emph{agent} qui orchestre les briques des
phases précédentes via une boucle \emph{Réfléchir-Agir-Observer} (ReAct).
On y aborde~: le décorateur \texttt{@tool} et la docstring comme
contrat avec le LLM, la construction de l'agent avec
\texttt{create\_agent} (LangChain 1.x), l'observation de la boucle ReAct
sur trois démos contrastées, et la gestion de la mémoire
conversationnelle pour un chat multi-tour. Chaque exercice suit la
structure habituelle (énoncé + objectifs en boîte, solution détaillée,
vérification LangSmith, récapitulatif). Une seconde partie
\emph{« Pour aller plus loin »} prolonge avec cinq exercices avancés sur
les agents~: ajouter un tool dynamique, observer un agent dans LangSmith,
paramétrer le \texttt{timeout} d'un tool, composer un tool avec un
sous-LLM, et streamer la réponse en deux modes.
\end{abstract}

\tableofcontents
\bigskip

\part{Exercices}

% #####################################################################
% #####################################################################
\section{Définir des \emph{tools} avec \texttt{@tool}}
% #####################################################################
% #####################################################################

\begin{tcolorbox}[title=Exercice 1. Quatre outils pour l'agent PyTutor]
Implémente quatre \emph{tools} que l'agent PyTutor pourra appeler~:
\texttt{search\_python\_docs} (recherche dans la doc Python via le RAG
de la Phase 2), \texttt{run\_python\_code} (exécute du code en
\texttt{subprocess}), \texttt{generate\_exercise} (utilise la chaîne
d'explication de la Phase 1), et \texttt{check\_pep8} (lance
\texttt{pycodestyle} sur du code).

Chaque tool doit être décoré avec \texttt{@tool} de
\texttt{langchain\_core.tools}, et sa \textbf{docstring} doit décrire
précisément à quoi il sert et \emph{quand} l'appeler — c'est ce que le
modèle va lire pour décider.

\textbf{Contrainte}~: gérer proprement les cas d'erreur (timeout du
subprocess, module absent, pas de résultat dans le RAG, valeur invalide
en argument) avec des messages qui parlent au LLM.

\bigskip
\noindent\textbf{Objectifs pédagogiques}
\begin{itemize}[leftmargin=*]
  \item Maîtriser le décorateur \texttt{@tool}~: il transforme une
        fonction Python en \emph{tool} appelable par un LLM.
  \item Comprendre que la \textbf{docstring} est envoyée au modèle comme
        description du tool — donc c'est du \emph{prompt engineering}
        appliqué à du code.
  \item Apprendre les patterns d'implémentation propres~: lazy init,
        \texttt{tempfile} avec nettoyage, \texttt{timeout}, fallbacks
        sur erreur.
  \item Réutiliser les briques des Phases 1 et 2 à l'intérieur d'un tool.
\end{itemize}
\end{tcolorbox}

\subsection{Solution}

\subsubsection{Architecture}

\begin{center}
\begin{tikzpicture}[node distance=5mm and 10mm]
  \node[llm] (model) {Modèle LLM};

  \node[tool, below left=14mm and 18mm of model]  (t1)
       {search\_python\_docs\\\scriptsize → RAG Chroma};
  \node[tool, below=14mm of model]                (t2)
       {run\_python\_code\\\scriptsize → subprocess};
  \node[tool, below right=14mm and 18mm of model] (t3)
       {generate\_exercise\\\scriptsize → chaîne LCEL};
  \node[tool, below=of t2]                        (t4)
       {check\_pep8\\\scriptsize → pycodestyle};

  \draw[flow] (model) -- (t1) node[label, midway, sloped, above]{appel};
  \draw[flow] (model) -- (t2);
  \draw[flow] (model) -- (t3);
  \draw[flow] (t2)    -- (t4);

  \draw[flow, dashed, draw=gray!60]  (t1.east) ..controls +(2,0) and +(2,0).. (model.east)
       node[label, midway, right]{résultat};
\end{tikzpicture}
\end{center}

Le LLM voit \textbf{4 tools déclarés} dans son prompt système (généré
par \texttt{create\_agent}). À chaque tour, il décide d'en appeler un
(ou plusieurs) en lisant leur docstring. Le résultat de chaque appel
retourne sous forme de \texttt{ToolMessage} dans l'historique, et le
LLM \emph{réfléchit} à nouveau.

\subsubsection{Le code complet}

\begin{pythoncode}
import subprocess
import sys
import tempfile
from pathlib import Path

from langchain_core.tools import tool

from phase1_explain import explain_chain_for_level
from phase2_rag import build_or_load_vectorstore


# Cache lazy : on charge le retriever une seule fois, à la première utilisation
_retriever = None

def _get_retriever():
    global _retriever
    if _retriever is None:
        vs = build_or_load_vectorstore()
        _retriever = vs.as_retriever(search_kwargs={"k": 3})
    return _retriever


@tool
def search_python_docs(query: str) -> str:
    """Recherche dans la documentation Python officielle.

    À utiliser pour toute question FACTUELLE sur Python : syntaxe d'une
    construction, comportement standard d'une fonction built-in, vocabulaire
    technique. PAS pour générer du code ou évaluer du code de l'élève.

    Args:
        query: La question ou les mots-clés à rechercher dans la documentation.
    """
    docs = _get_retriever().invoke(query)
    if not docs:
        return "Aucun résultat dans la doc Python pour cette recherche."
    return "\n\n---\n\n".join(
        f"Source: {d.metadata.get('source', '?')}\n{d.page_content}"
        for d in docs
    )


@tool
def run_python_code(code: str) -> str:
    """Exécute du code Python dans un sous-processus et renvoie stdout + stderr.

    À utiliser quand :
      - l'élève montre un code qui ne marche pas (pour voir la vraie erreur)
      - tu veux vérifier qu'un exemple que tu vas donner fonctionne réellement
      - tu veux mesurer la sortie d'un bout de code pour pouvoir l'expliquer

    Limites : timeout 5 secondes, environnement Python standard, pas de pip.

    Args:
        code: Le code Python à exécuter, complet et autonome.
    """
    with tempfile.NamedTemporaryFile(mode="w", suffix=".py", delete=False) as f:
        f.write(code)
        path = f.name
    try:
        result = subprocess.run(
            [sys.executable, path],
            capture_output=True, text=True, timeout=5,
        )
        parts = []
        if result.stdout:
            parts.append(f"STDOUT:\n{result.stdout}")
        if result.stderr:
            parts.append(f"STDERR:\n{result.stderr}")
        parts.append(f"Exit code: {result.returncode}")
        return "\n\n".join(parts)
    except subprocess.TimeoutExpired:
        return (
            "ERREUR: timeout après 5 secondes. Le code boucle probablement "
            "à l'infini (boucle while, récursion sans condition d'arrêt, etc.)"
        )
    finally:
        Path(path).unlink(missing_ok=True)


@tool
def generate_exercise(topic: str, difficulty: str = "intermédiaire") -> str:
    """Génère un exercice Python sur un concept avec ses questions de compréhension.

    À utiliser quand l'élève demande explicitement de la pratique, des
    exercices, ou pour vérifier sa compréhension d'un concept.

    Args:
        topic: Le concept Python à pratiquer (ex: "décorateurs", "list comprehensions").
        difficulty: Le niveau visé. Valeurs valides : "débutant", "intermédiaire", "avancé".
    """
    if difficulty not in {"débutant", "intermédiaire", "avancé"}:
        difficulty = "intermédiaire"

    chain = explain_chain_for_level(difficulty)
    explanation = chain.invoke({"concept": topic})

    questions = "\n".join(
        f"  Q{i}. {q.question}\n     -> Réponse attendue : {q.expected_answer}"
        for i, q in enumerate(explanation.comprehension_questions, 1)
    )
    return (
        f"Exercice sur '{topic}' (niveau {difficulty}) :\n\n"
        f"Code de référence :\n{explanation.code_example}\n\n"
        f"Questions :\n{questions}\n\n"
        f"Piège classique : {explanation.pitfall}"
    )


@tool
def check_pep8(code: str) -> str:
    """Vérifie la conformité PEP 8 d'un bout de code Python.

    À utiliser quand l'élève partage du code et veut savoir s'il est
    propre, ou quand tu donnes un exemple que tu veux valider stylistiquement.

    Args:
        code: Le code Python à analyser.
    """
    with tempfile.NamedTemporaryFile(mode="w", suffix=".py", delete=False) as f:
        f.write(code)
        path = f.name
    try:
        result = subprocess.run(
            [sys.executable, "-m", "pycodestyle", "--max-line-length=100", path],
            capture_output=True, text=True, timeout=5,
        )
        if not result.stdout.strip():
            return "Aucune violation PEP 8 détectée."
        cleaned = result.stdout.replace(path, "<code>")
        return f"Problèmes PEP 8 détectés :\n{cleaned}"
    except FileNotFoundError:
        return (
            "Le module pycodestyle n'est pas installé. "
            "Réponds à l'élève en lui suggérant `pip install pycodestyle`."
        )
    finally:
        Path(path).unlink(missing_ok=True)
\end{pythoncode}

\subsubsection{Explications pas à pas}

\subsubsection*{Étape 1 — Le décorateur \texttt{@tool}}

\texttt{@tool} de \texttt{langchain\_core.tools} transforme une fonction
Python en \texttt{BaseTool} consommable par un LLM. Ce qu'il fait sous
le capot~:

\begin{itemize}[leftmargin=*]
  \item Lit la \textbf{signature} (noms et types des arguments) pour
        générer un schéma JSON.
  \item Lit la \textbf{docstring} pour produire la \emph{description}
        du tool.
  \item Encapsule l'appel pour intercepter exceptions et timeouts (à
        compléter manuellement pour les cas avancés).
\end{itemize}

\begin{pythoncode}
@tool
def my_tool(query: str) -> str:
    """Description du tool, en français OK."""
    return "résultat"

# my_tool est maintenant un BaseTool, pas une fonction normale
my_tool.name           # "my_tool"
my_tool.description    # "Description du tool, en français OK."
my_tool.args_schema    # Pydantic model avec query: str
\end{pythoncode}

\subsubsection*{Étape 2 — La docstring est le contrat avec le LLM}

La docstring est \textbf{envoyée au modèle} comme description du tool.
Le modèle décide quand l'appeler en lisant cette description. Donc~:

\begin{itemize}[leftmargin=*]
  \item \textbf{À QUOI} sert le tool — première ligne, factuelle.
  \item \textbf{QUAND} l'utiliser — situations explicites.
  \item \textbf{QUAND PAS} l'utiliser — démarquage des autres tools.
  \item \textbf{Args} typés et décrits — pour que le modèle remplisse
        bien chaque paramètre.
\end{itemize}

\begin{quote}\itshape
Une docstring vague = un agent qui prend de mauvaises décisions. Pour
chaque tool, demande-toi : « si je donne CETTE docstring à un humain,
peut-il savoir QUAND l'appeler sans hésitation~?» Si non, précise.
\end{quote}

\subsubsection*{Étape 3 — Les paramètres typés}

Les types Python (\texttt{str}, \texttt{int}, \texttt{list[str]},
\texttt{Literal[\dots]}\dots) sont traduits en JSON Schema dans la
signature envoyée au LLM. Le modèle ne pourra pas appeler le tool avec
le mauvais type — c'est validé par Pydantic en amont.

Pour \texttt{generate\_exercise}, on accepte un \texttt{difficulty: str}
au lieu d'un \texttt{Literal} parce qu'on veut un fallback côté Python
plutôt qu'un crash si le modèle invente un niveau \texttt{"expert"}.
C'est un choix~: rigide via \texttt{Literal} pour échouer fort, ou
souple via \texttt{str + validation} pour dégrader proprement.

\subsubsection*{Étape 4 — Patterns d'implémentation}

\textbf{Lazy initialisation}. Le retriever du RAG (Phase 2) est lourd à
charger. On le \emph{cache} dans une variable de module avec
\texttt{global}~:

\begin{pythoncode}
_retriever = None

def _get_retriever():
    global _retriever
    if _retriever is None:
        vs = build_or_load_vectorstore()
        _retriever = vs.as_retriever(search_kwargs={"k": 3})
    return _retriever
\end{pythoncode}

Le retriever ne se charge que si le tool est effectivement appelé, et
seulement la première fois.

\textbf{Tempfile avec nettoyage}.
\texttt{tempfile.NamedTemporaryFile(delete=False)} + bloc
\texttt{try/finally} qui supprime explicitement~:

\begin{pythoncode}
with tempfile.NamedTemporaryFile(mode="w", suffix=".py", delete=False) as f:
    f.write(code)
    path = f.name
try:
    # ... utilise path ...
    pass
finally:
    Path(path).unlink(missing_ok=True)
\end{pythoncode}

\textbf{Timeout sur subprocess}. Une boucle \texttt{while True} dans le
code utilisateur bloquerait l'agent indéfiniment. Le \texttt{timeout=5}
de \texttt{subprocess.run} lève \texttt{TimeoutExpired}~:

\begin{pythoncode}
try:
    result = subprocess.run(..., timeout=5)
except subprocess.TimeoutExpired:
    return "ERREUR: timeout après 5 secondes. Le code boucle..."
\end{pythoncode}

\textbf{Messages d'erreur qui parlent au LLM}. Quand un tool échoue, le
message de retour devient l'observation que voit le modèle. Mauvais~:
\texttt{"Error 127"}. Bon~: \texttt{"Le module pycodestyle n'est pas
installé. Réponds à l'élève en lui suggérant pip install
pycodestyle."}~— le LLM sait alors quoi répondre.

\begin{quote}\itshape
Un tool n'a pas le droit de planter silencieusement, ni de lever une
exception non gérée~: dans les deux cas, l'agent perd la main et
l'utilisateur voit un crash. Toujours renvoyer un \emph{message clair}
en cas d'erreur, même un message d'erreur.
\end{quote}

\subsubsection*{Étape 5 — Réutiliser les briques des phases précédentes}

Chaque tool s'appuie sur ce qu'on a construit précédemment~:

\begin{center}
\begin{tabular}{p{4.5cm} p{10cm}}
\toprule
\textbf{Tool} & \textbf{Brique réutilisée} \\
\midrule
\texttt{search\_python\_docs} & Vector store + retriever (Phase 2). \\[1mm]
\texttt{run\_python\_code}    & \texttt{subprocess} standard, pas LangChain. \\[1mm]
\texttt{generate\_exercise}   & Chaîne \texttt{explain\_chain\_for\_level}
                                avec \texttt{with\_structured\_output} (Phase 1). \\[1mm]
\texttt{check\_pep8}          & \texttt{pycodestyle} en subprocess. \\
\bottomrule
\end{tabular}
\end{center}

\begin{quote}\itshape
Un tool peut contenir une chaîne LCEL complète, voire un autre RAG.
C'est le pattern « tool qui appelle un autre LLM en interne ». Très
puissant pour décomposer un système complexe en agents qui se passent
des tâches.
\end{quote}

\subsection{Vérification dans LangSmith}

Quand un tool est appelé par l'agent, LangSmith trace~:

\begin{center}
\begin{tikzpicture}[every node/.style={font=\footnotesize\ttfamily, anchor=west}]
  \node                              (a) {ToolNode (search\_python\_docs)};
  \node[below=1mm of a, xshift=4mm]  (b) {Input  : \{"query": "Comment fonctionne yield ?"\}};
  \node[below=1mm of b]              (c) {Output : "Source: docs.python.org/.../functional.html..."};
  \node[below=1mm of c, xshift=-4mm] (d) {ToolNode (run\_python\_code)};
  \node[below=1mm of d, xshift=4mm]  (e) {Input  : \{"code": "def f():..."\}};
  \node[below=1mm of e]              (f) {Output : "STDOUT: ... | Exit code: 0"};
\end{tikzpicture}
\end{center}

\begin{quote}\itshape
Si un tool ne s'appelle jamais alors qu'il \emph{devrait}~: regarde sa
docstring. Si un tool s'appelle \emph{trop}~: regarde aussi sa
docstring. C'est presque toujours là.
\end{quote}

\subsection*{Récapitulatif}
{\itshape\sloppy
\textit{Ce qu'il faut retenir :}
\begin{enumerate}[leftmargin=*]
  \item \texttt{@tool} transforme une fonction Python en outil
        consommable par un LLM. Signature et docstring deviennent un
        contrat envoyé au modèle.
  \item La docstring fait du \textbf{prompt engineering} sur le code~:
        explicite à \textbf{quoi} sert le tool, \textbf{quand} l'appeler,
        et \textbf{quand pas}.
  \item Patterns essentiels~: lazy init pour les ressources lourdes,
        \texttt{tempfile} avec nettoyage, \texttt{timeout} sur les
        subprocess, messages d'erreur qui parlent au LLM.
  \item Un tool peut contenir une chaîne LCEL, un RAG, ou n'importe
        quel autre traitement Python — c'est l'orchestration côté agent
        qui le rend pertinent.
\end{enumerate}
\par}

% #####################################################################
% #####################################################################
\section{Construire l'agent avec \texttt{create\_agent}}
% #####################################################################
% #####################################################################

\begin{tcolorbox}[title=Exercice 2. Assembler les tools dans un agent ReAct]
Construis un agent qui combine~: le modèle de chat, la liste des 4 tools
définis dans l'exercice 1, et un \texttt{system\_prompt} qui dicte le
comportement attendu (français, pédagogique, règles d'usage des tools).

\textbf{Contrainte}~: utiliser \texttt{create\_agent} de
\texttt{langchain.agents} (LangChain 1.x), qui renvoie un \emph{graphe
LangGraph} compilé prêt à \texttt{.invoke} ou \texttt{.stream}.

\bigskip
\noindent\textbf{Objectifs pédagogiques}
\begin{itemize}[leftmargin=*]
  \item Maîtriser \texttt{create\_agent(model, tools, system\_prompt)}~:
        la voie officielle LangChain 1.x pour les agents ReAct.
  \item Rédiger un \texttt{system\_prompt} qui couvre rôle, périmètre,
        règles de comportement, et choix de langue.
  \item Comprendre la signature d'invocation~:
        \texttt{agent.invoke(\{"messages": [\dots]\})}.
  \item Savoir que le retour de \texttt{create\_agent} est un
        \textbf{graphe LangGraph} (donc tout LangGraph s'applique).
\end{itemize}
\end{tcolorbox}

\subsection{Solution}

\subsubsection{Architecture}

\begin{center}
\begin{tikzpicture}[node distance=8mm and 14mm]
  \node[block] (model)  {Modèle LLM\\\scriptsize get\_model()};
  \node[block, right=of model] (tools) {Liste de 4 tools};
  \node[block, below=of model] (sysp)  {SYSTEM\_PROMPT\\\scriptsize règles + tools docs};

  \node[block, right=24mm of sysp, fill=teal!10, draw=teal!60!black,
        minimum width=44mm, minimum height=12mm] (ca) {create\_agent(\dots)};

  \node[store, below=of ca] (graph) {Graphe LangGraph\\compilé};

  \draw[flow] (model) -- (ca);
  \draw[flow] (tools) -- (ca);
  \draw[flow] (sysp)  -- (ca);
  \draw[flow] (ca)    -- (graph) node[label, midway, right]{retourne};
\end{tikzpicture}
\end{center}

\texttt{create\_agent} construit le graphe d'exécution de la boucle
ReAct~: un nœud LLM qui décide, un nœud \texttt{ToolNode} qui exécute,
et des arêtes conditionnelles qui les enchaînent jusqu'à la réponse
finale.

\subsubsection{Le code complet}

\begin{pythoncode}
from langchain.agents import create_agent
from config import get_model


TOOLS = [search_python_docs, run_python_code, generate_exercise, check_pep8]

SYSTEM_PROMPT = """Tu es PyTutor, un assistant pédagogique pour des élèves apprenant Python.

Tu as accès à 4 outils :
  - search_python_docs : pour récupérer du contenu de la doc Python officielle
  - run_python_code    : pour exécuter du code et observer ce qui se passe
  - generate_exercise  : pour générer un exercice et ses questions
  - check_pep8         : pour vérifier le style PEP 8

Règles de comportement :
1. Pour toute question factuelle sur Python (syntaxe, builtin, vocabulaire),
   utilise search_python_docs AVANT de répondre. Cite la source dans ta réponse.
2. Quand l'élève montre un code qui ne marche pas, exécute-le toi-même avec
   run_python_code AVANT d'expliquer pourquoi, pour voir la vraie erreur.
3. Quand tu donnes un exemple de code dans ta réponse, exécute-le d'abord pour
   être sûr qu'il marche.
4. Sois pédagogique : explique POURQUOI, pas seulement COMMENT.
5. Adapte ton vocabulaire au niveau de l'élève (déduis-le de ses questions).
6. Réponds en français.
"""


def build_agent():
    """Construit l'agent PyTutor avec ses 4 tools."""
    return create_agent(
        model=get_model(),
        tools=TOOLS,
        system_prompt=SYSTEM_PROMPT,
    )
\end{pythoncode}

\subsubsection{Explications pas à pas}

\subsubsection*{Étape 1 — La signature de \texttt{create\_agent}}

\begin{pythoncode}
create_agent(
    model,                # instance ChatModel ou string "provider:model"
    tools,                # liste de @tool callables
    system_prompt=None,   # str, remplace l'ancien `prompt`
    # ... d'autres options avancées
)
\end{pythoncode}

Trois paramètres essentiels~:

\begin{itemize}[leftmargin=*]
  \item \textbf{\texttt{model}}~: une instance issue de \texttt{get\_model()}
        (ou n'importe quel \texttt{BaseChatModel}). Ou une string courte
        \texttt{"anthropic:claude-3-5-sonnet-20241022"}.
  \item \textbf{\texttt{tools}}~: la liste de fonctions décorées
        \texttt{@tool}. L'ordre n'a pas d'importance fonctionnelle,
        mais peut influencer marginalement les décisions du modèle
        (priorité au premier tool listé en cas d'ambiguïté).
  \item \textbf{\texttt{system\_prompt}}~: instructions générales. C'est
        le \emph{contrat global} de l'agent — son rôle, ses règles, son
        ton.
\end{itemize}

\subsubsection*{Étape 2 — Anatomie d'un \texttt{system\_prompt} efficace}

Le \texttt{SYSTEM\_PROMPT} de PyTutor a quatre sections distinctes~:

\begin{enumerate}[leftmargin=*]
  \item \textbf{Identité} (« Tu es PyTutor, un assistant pédagogique
        \dots »).
  \item \textbf{Inventaire des tools} (la même liste qu'aux yeux du
        modèle, en clair). Redondant mais utile~: ancre la décision.
  \item \textbf{Règles d'usage} numérotées qui couvrent les cas
        critiques (utilise X avant Y, exécute le code avant de
        l'expliquer, etc.).
  \item \textbf{Méta-comportement} (langue, ton, public).
\end{enumerate}

\begin{quote}\itshape
Le \texttt{system\_prompt} n'est pas un texte décoratif~: c'est le
\textbf{principal levier de qualité} d'un agent. Une règle bien
formulée évite des centaines d'appels de tool inutiles. Investis-y du
temps.
\end{quote}

\subsubsection*{Étape 3 — Le graphe LangGraph sous le capot}

\texttt{create\_agent} retourne un \textbf{graphe LangGraph compilé}, pas
un simple \texttt{Runnable}. Visuellement, ce graphe a la forme~:

\begin{center}
\begin{tikzpicture}[node distance=5mm and 18mm]
  \node[llm] (a) {agent (LLM)};
  \node[tool, right=of a] (t) {tools (ToolNode)};
  \node[block, below=of a] (end) {END};

  \draw[flow] (a) -- (t) node[label, midway, above]{tool\_calls non vide};
  \draw[flow] (t) -- (a) node[label, midway, below]{ToolMessage};
  \draw[flow] (a) -- (end) node[label, midway, right]{pas de tool\_calls};
\end{tikzpicture}
\end{center}

À chaque tour, le nœud LLM examine la conversation. Soit il décide
d'appeler des tools (le graphe passe au \texttt{ToolNode}, exécute les
appels, revient au LLM), soit il produit une réponse finale (le graphe
se termine à \texttt{END}).

\subsubsection*{Étape 4 — Invocation : \texttt{\{"messages": [\dots]\}}}

L'agent s'utilise avec un dict contenant la clé \texttt{messages}~:

\begin{pythoncode}
from langchain_core.messages import HumanMessage

agent = build_agent()
result = agent.invoke({
    "messages": [HumanMessage(content="Comment fonctionne yield ?")]
})

# result["messages"] contient TOUT l'historique :
# [SystemMessage, HumanMessage, AIMessage(tool_calls=...), ToolMessage, AIMessage]
\end{pythoncode}

Le \texttt{system\_prompt} est injecté automatiquement au début par
\texttt{create\_agent}~: tu n'as pas à le passer toi-même dans
\texttt{messages}.

\subsubsection*{Étape 5 — Le retour est l'historique complet}

\texttt{result["messages"]} contient \textbf{toute la conversation},
pas seulement la dernière réponse~:

\begin{itemize}[leftmargin=*]
  \item Le \texttt{SystemMessage} (ajouté par \texttt{create\_agent}).
  \item Le \texttt{HumanMessage} initial.
  \item Les \texttt{AIMessage} intermédiaires, parfois avec
        \texttt{tool\_calls}.
  \item Les \texttt{ToolMessage} qui contiennent les résultats des tools.
  \item Le \texttt{AIMessage} final (la réponse à l'utilisateur).
\end{itemize}

C'est cette structure qu'on va lire dans les démos de l'exercice 3.

\subsection{Vérification dans LangSmith}

\subsubsection*{L'arbre d'exécution attendu}

\begin{center}
\begin{tikzpicture}[every node/.style={font=\footnotesize\ttfamily, anchor=west}]
  \node                              (a) {LangGraph (agent compilé)};
  \node[below=1mm of a, xshift=4mm]  (b) {agent (LLM) ← tour 1};
  \node[below=1mm of b]              (c) {tools (ToolNode) ← exécute le(s) appel(s)};
  \node[below=1mm of c]              (d) {agent (LLM) ← tour 2};
  \node[below=1mm of d]              (e) {(éventuellement tools $\to$ agent à nouveau)};
  \node[below=1mm of e]              (f) {END};
\end{tikzpicture}
\end{center}

Trois choses à regarder dans une trace d'agent~:

\begin{enumerate}[leftmargin=*]
  \item \textbf{Le \texttt{system\_prompt} effectif} dans le premier
        appel LLM (onglet \emph{Input}). Vérifie qu'il contient tes
        règles ET les descriptions des tools.
  \item \textbf{Les \texttt{tool\_calls}} dans chaque \texttt{AIMessage}~:
        nom du tool appelé, arguments choisis par le LLM.
  \item \textbf{Les \texttt{ToolMessage}}~: ce que les tools ont
        \emph{vraiment} renvoyé au LLM. Une réponse incorrecte vient
        souvent d'un tool qui a renvoyé un truc inattendu.
\end{enumerate}

\subsection*{Récapitulatif}
{\itshape\sloppy
\textit{Ce qu'il faut retenir :}
\begin{enumerate}[leftmargin=*]
  \item \texttt{create\_agent(model, tools, system\_prompt)} construit
        un agent ReAct prêt à l'emploi en LangChain 1.x.
  \item Le retour est un \textbf{graphe LangGraph compilé}~: tout
        LangGraph s'applique (state, persistence, streaming, etc.).
  \item Le \texttt{system\_prompt} doit avoir 4 parties~: identité,
        inventaire des tools, règles numérotées, méta-comportement.
  \item Invocation~:
        \texttt{agent.invoke(\{"messages": [HumanMessage(\dots)]\})}.
  \item Le retour est l'historique \textbf{complet} de la conversation
        (System / Human / AI / Tool / AI / \dots).
\end{enumerate}
\par}

% #####################################################################
% #####################################################################
\section{Observer la boucle Réfléchir-Agir-Observer (ReAct)}
% #####################################################################
% #####################################################################

\begin{tcolorbox}[title=Exercice 3. Trois démos contrastées pour voir la boucle ReAct]
Lance l'agent sur trois requêtes qui déclenchent des comportements
différents~:

\begin{enumerate}[leftmargin=*]
  \item \textbf{Question factuelle}~: doit déclencher
        \texttt{search\_python\_docs}.
  \item \textbf{Code qui boucle}~: doit déclencher
        \texttt{run\_python\_code} (l'agent voit le timeout en clair).
  \item \textbf{Question multi-étape}~: doit déclencher
        \emph{plusieurs} tools en chaîne (recherche + exercice).
\end{enumerate}

Écris une fonction \texttt{print\_trace(messages)} qui formate
l'historique de la conversation pour qu'on \emph{voie} le déroulé de
la boucle ReAct.

\textbf{Contrainte}~: distinguer visuellement \texttt{HumanMessage},
\texttt{AIMessage} (texte vs \texttt{tool\_calls}), et
\texttt{ToolMessage} (résultat des tools).

\bigskip
\noindent\textbf{Objectifs pédagogiques}
\begin{itemize}[leftmargin=*]
  \item Lire l'historique de messages d'un agent~: ordre, types, contenu.
  \item Comprendre la boucle \textbf{Réfléchir-Agir-Observer}~: le LLM
        décide $\to$ le tool s'exécute $\to$ le LLM observe $\to$ il
        re-décide jusqu'à pouvoir répondre.
  \item Identifier quand un agent décide de \emph{s'arrêter}~:
        l'\texttt{AIMessage} sans \texttt{tool\_calls} est la réponse
        finale.
  \item Construire un affichage lisible pour le déroulé pédagogique.
\end{itemize}
\end{tcolorbox}

\subsection{Solution}

\subsubsection{Architecture}

\begin{center}
\begin{tikzpicture}[node distance=4mm and 10mm]
  \node[block]                          (start) {HumanMessage};
  \node[llm,   below=of start]          (m1)    {Modèle\\\scriptsize tour 1 (Réfléchir)};
  \node[tool,  below=of m1]             (t1)    {tool\_call exécuté\\\scriptsize ToolMessage en sortie (Observer)};
  \node[llm,   below=of t1]             (m2)    {Modèle\\\scriptsize tour 2 (re-Réfléchir)};
  \node[block, below=of m2]             (end)   {AIMessage final\\\scriptsize (pas de tool\_call)};

  \draw[flow] (start) -- (m1);
  \draw[flow] (m1)    -- (t1) node[label, midway, right]{Agir};
  \draw[flow] (t1)    -- (m2);
  \draw[flow] (m2)    -- (end);

  \draw[flow, dashed, draw=orange!70!black]
       (m2.west) ..controls +(-2,0) and +(-2,0).. (t1.west)
       node[label, midway, left]{loop si tool\_calls};
\end{tikzpicture}
\end{center}

À chaque tour~:
\begin{enumerate}[leftmargin=*]
  \item Le LLM \emph{réfléchit}~: examine la conversation et décide.
  \item S'il décide d'appeler un ou plusieurs tools, il les liste dans
        \texttt{AIMessage.tool\_calls}.
  \item Le \texttt{ToolNode} \emph{agit}~: exécute chaque appel et
        produit un \texttt{ToolMessage} par tool.
  \item Le LLM \emph{observe}~: ces \texttt{ToolMessage} entrent dans
        son contexte au tour suivant.
  \item Boucle jusqu'à ce que le LLM produise un \texttt{AIMessage}
        \textbf{sans \texttt{tool\_calls}}~— c'est la réponse finale.
\end{enumerate}

\subsubsection{Le code complet}

\begin{pythoncode}
from langchain_core.messages import (
    HumanMessage, AIMessage, ToolMessage, SystemMessage,
)


def print_trace(messages, from_index: int = 0) -> None:
    """Affiche le déroulé de la boucle ReAct depuis un index donné."""
    for msg in messages[from_index:]:
        if isinstance(msg, SystemMessage):
            continue   # on n'affiche pas le system prompt
        elif isinstance(msg, HumanMessage):
            print(f"\n[USER]:\n  {msg.content}")
        elif isinstance(msg, AIMessage):
            if msg.content:
                print(f"\n[AI]:\n  {msg.content}")
            if hasattr(msg, "tool_calls") and msg.tool_calls:
                print(f"\n[AI] décide d'appeler {len(msg.tool_calls)} tool(s) :")
                for tc in msg.tool_calls:
                    args_preview = str(tc["args"])[:100]
                    suffix = "..." if len(str(tc["args"])) > 100 else ""
                    print(f"   -> {tc['name']}({args_preview}{suffix})")
        elif isinstance(msg, ToolMessage):
            preview = msg.content[:300].replace("\n", "\n     ")
            ellipsis = "..." if len(msg.content) > 300 else ""
            print(f"\n[TOOL {msg.name}]:\n     {preview}{ellipsis}")


def demo_factual_question() -> None:
    """Question factuelle : devrait déclencher search_python_docs."""
    agent = build_agent()
    result = agent.invoke({
        "messages": [HumanMessage(content="À quoi sert exactement le mot-clé `yield` ?")]
    })
    print_trace(result["messages"])


def demo_buggy_code() -> None:
    """Code qui boucle à l'infini : devrait déclencher run_python_code."""
    buggy = """\
def double_each(numbers):
    for n in numbers:
        numbers.append(n * 2)
    return numbers

print(double_each([1, 2, 3]))
"""
    agent = build_agent()
    result = agent.invoke({
        "messages": [HumanMessage(content=(
            "Mon code ne se termine jamais. Tu peux m'aider ?\n\n"
            f"```python\n{buggy}```"
        ))]
    })
    print_trace(result["messages"])


def demo_multi_tool() -> None:
    """Question complexe : devrait chaîner plusieurs tools."""
    agent = build_agent()
    result = agent.invoke({
        "messages": [HumanMessage(content=(
            "Explique-moi les list comprehensions en t'appuyant sur la doc, "
            "puis génère-moi un exercice de niveau intermédiaire pour pratiquer."
        ))]
    })
    print_trace(result["messages"])
\end{pythoncode}

\subsubsection{Explications pas à pas}

\subsubsection*{Étape 1 — La structure de l'historique de messages}

\texttt{result["messages"]} est une liste hétérogène, dans l'ordre
chronologique~:

\begin{center}
\begin{tabular}{p{3.5cm} p{11cm}}
\toprule
\textbf{Type} & \textbf{Rôle} \\
\midrule
\texttt{SystemMessage}
& Le \texttt{system\_prompt} injecté par \texttt{create\_agent}. Toujours
en première position. \\[1mm]
\texttt{HumanMessage}
& Une requête utilisateur. \\[1mm]
\texttt{AIMessage}
& Une réponse (partielle ou finale) du LLM. Peut contenir~:
\texttt{content} (texte), \texttt{tool\_calls} (liste d'appels à
exécuter), ou les deux. \\[1mm]
\texttt{ToolMessage}
& Le résultat d'\emph{un} appel de tool. Contient \texttt{name} (nom du
tool), \texttt{content} (résultat str), \texttt{tool\_call\_id}
(corrélation avec l'\texttt{AIMessage.tool\_calls}). \\
\bottomrule
\end{tabular}
\end{center}

\subsubsection*{Étape 2 — La boucle ReAct en pratique}

Pour la démo \emph{question factuelle}, on observe typiquement~:

\begin{enumerate}[leftmargin=*]
  \item \texttt{HumanMessage}~: « À quoi sert yield~? »
  \item \texttt{AIMessage} (tour 1, \textbf{Réfléchir})~: pas de
        \texttt{content}, juste \texttt{tool\_calls} =
        \texttt{[search\_python\_docs(query="yield")]}.
  \item \texttt{ToolMessage} (\textbf{Agir + Observer})~: contient les
        chunks de la doc sur \texttt{yield}.
  \item \texttt{AIMessage} (tour 2, \textbf{re-Réfléchir})~:
        \texttt{content} = réponse finale, pas de \texttt{tool\_calls}.
\end{enumerate}

C'est une boucle ReAct à \textbf{une seule itération}~: l'agent réfléchit,
agit une fois, observe, puis répond.

Pour la démo \emph{multi-tool}, la boucle a typiquement 2 itérations~:
l'agent appelle d'abord \texttt{search\_python\_docs} pour la doc,
\emph{observe} le résultat, puis appelle \texttt{generate\_exercise}
pour l'exercice, observe à nouveau, puis répond.

\subsubsection*{Étape 3 — La condition d'arrêt}

L'agent s'arrête \textbf{exactement} quand le LLM produit un
\texttt{AIMessage} sans \texttt{tool\_calls}. Ce n'est pas un nombre
maximal de tours imposé par défaut~: c'est le LLM qui décide qu'il a
assez d'informations pour répondre.

\begin{quote}\itshape
\textbf{Garde-fou de sécurité} : par défaut, le graphe a un
\texttt{recursion\_limit} (typiquement 25 tours) qui force l'arrêt si
l'agent boucle sans converger. C'est rare en pratique mais utile pour
détecter un bug de prompt (« utilise toujours \texttt{search\_python\_docs}~»
qui crée une boucle infinie de recherches).
\end{quote}

\subsubsection*{Étape 4 — \texttt{print\_trace} : afficher proprement}

Trois choix d'affichage~:

\begin{itemize}[leftmargin=*]
  \item Distinguer \textbf{visuellement} chaque type de message avec un
        préfixe (\texttt{[USER]}, \texttt{[AI]}, \texttt{[TOOL nom]}).
        En production, on utiliserait des emojis ou des couleurs.
  \item \textbf{Tronquer} les contenus longs (\texttt{[:300]}) pour
        garder un déroulé lisible. Sans troncature, un \texttt{ToolMessage}
        avec 4000 caractères de doc rend la trace illisible.
  \item \textbf{Indenter} les contenus pour qu'on voit la structure
        d'un coup d'œil.
\end{itemize}

\begin{tcolorbox}[
  colback=white,
  colframe=black!50,
  coltitle=black,
  fonttitle=\footnotesize\bfseries,
  fontupper=\footnotesize,
  title=Sortie de l'exécution,
  breakable,
]
\begin{verbatim}
======================================================================
DEMO 1 -- Question factuelle
======================================================================
[OK] Rechargement de l'index depuis ./chroma_db

[USER]:
  A quoi sert exactement le mot-clé `yield` ?

[AI]:
  Je vais chercher la documentation officielle pour te donner une
  explication précise.
  -> tool_use: search_python_docs({'query': 'yield keyword
                                              generator'})

[TOOL search_python_docs]:
     Source: https://docs.python.org/3/howto/functional.html
     are messy. In Python 2.5 there's a simple way to pass values
     into a generator. yield became an expression, returning a
     value that can be assigned to a variable or otherwise operated
     on: val = (yield i) ...

[AI]:
  Parfait ! Voici l'explication :

  ## Le mot-clé `yield`
  **`yield` sert à créer des générateurs** -- des fonctions qui
  produisent une séquence de valeurs **une à la fois**, sans tout
  charger en mémoire d'un coup.

  ### Différence clé avec `return`
  - **`return`** : termine la fonction et retourne une valeur
  - **`yield`** : **suspend** la fonction, retourne une valeur, et
    **mémorise l'état** pour reprendre plus tard

  -> tool_use: run_python_code({...comparaison classique/yield...})

[TOOL run_python_code]:
     STDOUT:
     === Approche classique ===
     Type : <class 'list'>
     Contenu : [0, 1, 2, 3, 4]

     === Approche avec yield ===
     Type : <class 'generator'>
     Contenu : <generator object nombres_generateur at 0x1029af780>
     Valeurs du générateur :
     0
     1
     2
     3
     4
     Exit code: 0

[AI]:
  ### Avantages des générateurs
  1. **Economie de mémoire** : les valeurs sont générées à la
     demande, pas stockées
  2. **Paresse** (lazy evaluation) : utile pour les grandes
     séquences ou infinies
  3. **Lisibilité** : code plus simple et naturel

  -> tool_use: run_python_code({...comparaison taille mémoire...})

[TOOL run_python_code]:
     STDOUT:
     Taille d'une liste de 1M nombres : 8000056 bytes
     Taille d'un générateur : 192 bytes
     Exit code: 0

[AI]:
  Enorme différence !

  ### Résumé
  | Aspect      | `return`         | `yield`            |
  |-------------|------------------|--------------------|
  | Retourne    | Une valeur unique| Un générateur      |
  | Mémoire     | Tout en mémoire  | A la demande       |
  | Reprendre ? | Non, terminée    | Oui, état préservé |
  | Cas d'usage | Résultat unique  | Séquences          |

  Source : Documentation Python officielle - Generators

======================================================================
DEMO 2 -- Debugging d'un code qui boucle
======================================================================

[USER]:
  Mon code ne se termine jamais. Tu peux m'aider ?

  ```python
  def double_each(numbers):
      for n in numbers:
          numbers.append(n * 2)
      return numbers

  print(double_each([1, 2, 3]))
  ```

[AI]:
  Je vais exécuter ton code pour voir ce qui se passe :
  -> tool_use: run_python_code({'code': '...'})

[TOOL run_python_code]:
     ERREUR: timeout après 5 secondes. Le code boucle probablement
     à l'infini (boucle while, génération récursive sans condition
     d'arrêt, etc.)

[AI]:
  Exactement ! Ton code crée une **boucle infinie**. Voici pourquoi :

  ## Le problème
  Tu modifies la liste **pendant que tu la parcours** :
    Iteration 1 : n=1, append(2)   -> [1, 2, 3, 2]
    Iteration 2 : n=2, append(4)   -> [1, 2, 3, 2, 4]
    Iteration 3 : n=3, append(6)   -> [1, 2, 3, 2, 4, 6]
    Iteration 4 : n=2, append(4)   -> [1, 2, 3, 2, 4, 6, 4]  <- continue !

  ## La solution
  **Option 1 : Créer une nouvelle liste**
  ```python
  def double_each(numbers):
      result = []
      for n in numbers:
          result.append(n * 2)
      return result
  ```

  **Option 2 : Parcourir une copie**
  ```python
  def double_each(numbers):
      for n in numbers[:]:  # [:] crée une copie
          numbers.append(n * 2)
      return numbers
  ```

  **Option 3 : List comprehension (plus pythonique)**
  ```python
  def double_each(numbers):
      return [n * 2 for n in numbers]
  ```

======================================================================
DEMO 3 -- Multi-tool (explication + exercice)
======================================================================

[USER]:
  Explique-moi les list comprehensions en t'appuyant sur la doc,
  puis génère-moi un exercice de niveau intermédiaire pour pratiquer.

[AI]:
  Je vais d'abord consulter la documentation officielle sur les
  list comprehensions, puis je te générerai un exercice.
  -> tool_use: search_python_docs({'query': 'list comprehensions'})
  -> tool_use: generate_exercise({'topic': 'list comprehensions',
                                  'difficulty': 'intermédiaire'})

[TOOL search_python_docs]:
     Source: https://docs.python.org/3/tutorial/datastructures.html
     For example, assume we want to create a list of squares, like:
     >>> squares = []
     >>> for x in range(10):
     ...     squares.append(x ** 2)
     >>> squares
     [0, 1, 4, 9, 16, 25, 36, 49, 64, 81]
     Note that this creates (or overwrites) a vari...

[TOOL generate_exercise]:
     Exercice sur 'list comprehensions' (niveau intermédiaire) :
     Code de référence :
     # Approche classique avec une boucle
     nombres = [1, 2, 3, 4, 5]
     carres_classique = []
     for n in nombres:
         carres_classique.append(n ** 2)
     # Approche avec list comprehension
     carres = [n ** 2 for n in...

[AI]:
  ## Les List Comprehensions expliquées

  Une **list comprehension** est une syntaxe compacte pour créer une
  liste. Selon la doc, elle se compose de :
  - Des **crochets** `[ ]`
  - Une **expression** à évaluer
  - Une ou plusieurs **clauses `for`**
  - Zéro ou plusieurs **clauses `if`** (optionnelles)

  **Syntaxe générale :**
  ```python
  [expression for variable in iterable if condition]
  ```

  ### Exemples clés
  **1. Avec une condition (filtrage) :**
  ```python
  [x for x in range(10) if x % 2 == 0]   # [0, 2, 4, 6, 8]
  ```

  **2. Avec plusieurs boucles (produit cartésien) :**
  ```python
  [(x, y) for x in [1, 2, 3] for y in [3, 1, 4] if x != y]
  ```

  ## Exercice de niveau intermédiaire
  **Q1.** Qu'est-ce qu'une list comprehension et quel est son
  avantage par rapport à une boucle for classique ?

  **Q2.** Ecris une list comprehension qui crée une liste contenant
  les carrés des nombres pairs de 1 à 10.

  **Q3.** Pourquoi serait-il mauvais d'utiliser une list comprehension
  pour une transformation très complexe impliquant plusieurs
  conditions imbriquées et appels de fonctions ?

  Piège classique à éviter :
   - Confondre l'ordre : `[for x in liste expression]` (FAUX)
   - Bon ordre        : `[expression for x in liste]` (CORRECT)
\end{verbatim}
\end{tcolorbox}

\subsection{Vérification dans LangSmith}

LangSmith trace chaque agent run avec une visualisation arborescente~:

\begin{center}
\begin{tikzpicture}[every node/.style={font=\footnotesize\ttfamily, anchor=west}]
  \node                              (a) {LangGraph run};
  \node[below=1mm of a, xshift=4mm]  (b) {agent (LLM, tour 1)};
  \node[below=1mm of b, xshift=4mm]  (c) {Input  : [System, Human]};
  \node[below=1mm of c]              (d) {Output : tool\_calls = [\{search\_python\_docs\}]};
  \node[below=1mm of d, xshift=-4mm] (e) {tools (ToolNode)};
  \node[below=1mm of e, xshift=4mm]  (f) {search\_python\_docs(query="yield")};
  \node[below=1mm of f]              (g) {  Output : "Source: ..."};
  \node[below=1mm of g, xshift=-4mm] (h) {agent (LLM, tour 2)};
  \node[below=1mm of h, xshift=4mm]  (i) {Input  : [System, Human, AI, Tool]};
  \node[below=1mm of i]              (j) {Output : "Le mot-clé yield..."};
\end{tikzpicture}
\end{center}

\begin{quote}\itshape
La vue \emph{timeline} dans LangSmith montre la durée de chaque tour~:
on voit instantanément si un tool est lent (souvent
\texttt{search\_python\_docs} avec embedding + recherche), si le LLM est
lent (gros prompt), ou si la boucle a \emph{trop} d'itérations
(probablement un bug de prompt).
\end{quote}

\subsection*{Récapitulatif}
{\itshape\sloppy
\textit{Ce qu'il faut retenir :}
\begin{enumerate}[leftmargin=*]
  \item L'historique d'un agent est une liste hétérogène~:
        \texttt{System}, \texttt{Human}, \texttt{AIMessage} (avec ou
        sans \texttt{tool\_calls}), \texttt{ToolMessage}.
  \item Boucle ReAct~: \textbf{Réfléchir} (LLM décide) $\to$
        \textbf{Agir} (tools exécutent) $\to$ \textbf{Observer} (LLM
        relit) $\to$ re-Réfléchir, jusqu'à un \texttt{AIMessage} sans
        \texttt{tool\_calls}.
  \item L'agent décide seul de s'arrêter~: il n'y a pas de nombre fixe
        de tours, juste un garde-fou \texttt{recursion\_limit} en cas de
        bug.
  \item Pour debugger un agent~: \textbf{lit l'historique}, regarde les
        \texttt{tool\_calls} et les \texttt{ToolMessage}. La plupart des
        problèmes sont là.
\end{enumerate}
\par}

% #####################################################################
% #####################################################################
\section{Mémoire conversationnelle multi-tours}
% #####################################################################
% #####################################################################

\begin{tcolorbox}[title=Exercice 4. Un chat où l'agent se souvient des tours précédents]
Construis une boucle de chat (\emph{REPL}) où l'utilisateur peut poser
plusieurs questions successives, et où l'agent \textbf{se souvient} des
tours précédents (utile pour les « peux-tu approfondir~? », « et avec
un exemple~? »\dots).

\textbf{Contrainte}~: \texttt{create\_agent} ne stocke \emph{rien} entre
les invocations. La mémoire est à maintenir côté client en accumulant
l'historique de \texttt{messages} et en le réinjectant à chaque tour.

\bigskip
\noindent\textbf{Objectifs pédagogiques}
\begin{itemize}[leftmargin=*]
  \item Comprendre que l'agent par défaut est \textbf{stateless}~: chaque
        \texttt{invoke} est indépendant.
  \item Construire la mémoire côté client en accumulant la liste de
        messages.
  \item Quantifier le coût (en tokens) d'une mémoire qui s'allonge à
        chaque tour.
  \item Connaître les alternatives (\texttt{checkpointer} LangGraph) pour
        une persistence native côté serveur.
\end{itemize}
\end{tcolorbox}

\subsection{Solution}

\subsubsection{Architecture}

\begin{center}
\begin{tikzpicture}[node distance=6mm and 14mm]
  \node[block] (build) {build\_agent()\\\scriptsize une seule fois};
  \node[data, below=of build] (hist) {history = []\\\scriptsize liste des messages};

  \node[block, below=of hist] (loop) {while True:};
  \node[block, below=of loop, fill=teal!10] (read) {user\_input = input(">")};
  \node[block, below=of read] (append) {history.append(HumanMessage(\dots))};
  \node[block, below=of append] (invoke) {result = agent.invoke(\{"messages": history\})};
  \node[block, below=of invoke] (replace) {history = result["messages"]\\\scriptsize remplace par l'historique enrichi};
  \node[block, below=of replace] (print) {print\_trace(history, from\_index=\dots)};

  \draw[flow] (build)   -- (hist);
  \draw[flow] (hist)    -- (loop);
  \draw[flow] (loop)    -- (read);
  \draw[flow] (read)    -- (append);
  \draw[flow] (append)  -- (invoke);
  \draw[flow] (invoke)  -- (replace);
  \draw[flow] (replace) -- (print);

  \draw[flow, dashed, draw=orange!70!black]
       (print.east) ..controls +(2,0) and +(2,0).. (read.east)
       node[label, midway, right]{tour suivant};
\end{tikzpicture}
\end{center}

\textbf{Important}~: l'agent est construit une seule fois (hors boucle).
À chaque tour, on accumule un nouveau \texttt{HumanMessage}, on
\texttt{invoke} avec \emph{tout} l'historique, et on remplace
l'historique par celui que renvoie l'agent (qui contient maintenant les
nouveaux \texttt{AIMessage} et \texttt{ToolMessage}).

\subsubsection{Le code complet}

\begin{pythoncode}
def chat_loop() -> None:
    """Boucle de chat interactive avec mémoire conversationnelle.

    create_agent ne stocke RIEN entre les invocations. Pour avoir une
    conversation continue, on accumule les messages côté client et on les
    renvoie à chaque tour.
    """
    print("MODE CHAT — Ctrl+C ou ligne vide pour quitter")

    agent = build_agent()
    history: list = []   # accumulera HumanMessage / AIMessage / ToolMessage

    while True:
        try:
            user_input = input("\n> ").strip()
        except (EOFError, KeyboardInterrupt):
            print("\nÀ bientôt.")
            return

        if not user_input:
            print("Au revoir.")
            return

        # 1. On ajoute le nouveau message au contexte complet
        history.append(HumanMessage(content=user_input))
        index_before = len(history) - 1

        # 2. L'agent renvoie l'historique complet enrichi
        result  = agent.invoke({"messages": history})
        history = result["messages"]

        # 3. On n'affiche que les messages produits ce tour-ci
        print_trace(history, from_index=index_before)
\end{pythoncode}

\subsubsection{Explications pas à pas}

\subsubsection*{Étape 1 — Pourquoi l'agent est stateless par défaut}

\texttt{create\_agent} retourne un graphe LangGraph compilé qui, par
défaut, n'a \textbf{pas de checkpointer}. Chaque \texttt{.invoke()}
exécute le graphe \emph{from scratch} avec uniquement les messages
qu'on lui donne en entrée. Aucune mémoire entre appels.

C'est un choix de design~: la persistence est un \emph{problème
séparé} (où stocker, comment indexer, comment expirer). LangChain ne
veut pas imposer un backend particulier.

\subsubsection*{Étape 2 — Maintenir l'historique côté client}

Le pattern est simple~: on garde une liste \texttt{history}, on y ajoute
le nouveau message utilisateur, on \texttt{invoke}, et on remplace
\texttt{history} par ce que l'agent renvoie~:

\begin{pythoncode}
history.append(HumanMessage(content=user_input))
result  = agent.invoke({"messages": history})
history = result["messages"]
\end{pythoncode}

Après \texttt{invoke}, \texttt{history} contient~:
\begin{itemize}[leftmargin=*]
  \item Tout l'historique précédent.
  \item Le nouveau \texttt{HumanMessage}.
  \item Les éventuels \texttt{AIMessage(tool\_calls=\dots)} intermédiaires.
  \item Les éventuels \texttt{ToolMessage} avec leurs résultats.
  \item Le \texttt{AIMessage} final.
\end{itemize}

Au tour suivant, ces messages sont \emph{tous} renvoyés à l'agent — le
LLM peut donc se référer aux échanges précédents.

\subsubsection*{Étape 3 — Le coût en tokens augmente à chaque tour}

\begin{center}
\begin{tabular}{lll}
\toprule
\textbf{Tour} & \textbf{Tokens d'entrée du LLM} & \textbf{Cumul} \\
\midrule
1   & system + question 1 ($\sim$300 tokens) & 300 \\
2   & + AI 1 + Human 2 ($\sim$800 tokens)    & 1100 \\
3   & + AI 2 + Human 3 ($\sim$1500 tokens)   & 2600 \\
\dots & \dots & \dots \\
10  & cumul $\sim$15000 tokens               & 15000 \\
\bottomrule
\end{tabular}
\end{center}

À chaque tour, on \emph{paie} l'intégralité de l'historique précédent en
tokens d'entrée. Sur des conversations longues, ça devient coûteux.

\begin{quote}\itshape
\textbf{Mitigation}~: tronquer l'historique (garder les N derniers
messages), ou \emph{résumer} les anciens messages avec une chaîne LCEL
quand la conversation dépasse une certaine taille (technique de la
\emph{summary memory}).
\end{quote}

\subsubsection*{Étape 4 — Termination propre}

Trois façons de sortir du chat~:

\begin{itemize}[leftmargin=*]
  \item \texttt{Ctrl+C}~: \texttt{input()} lève \texttt{KeyboardInterrupt}.
  \item \texttt{Ctrl+D}~: \texttt{input()} lève \texttt{EOFError}.
  \item \textbf{Ligne vide}~: \texttt{input()} renvoie \texttt{""}, on
        traite explicitement.
\end{itemize}

\subsubsection*{Étape 5 — Alternative serveur : le \texttt{checkpointer} LangGraph}

Pour une persistence native (utile dans un backend web), LangGraph
fournit des \emph{checkpointers}~:

\begin{pythoncode}
from langgraph.checkpoint.memory import MemorySaver

agent = create_agent(
    model=get_model(),
    tools=TOOLS,
    system_prompt=SYSTEM_PROMPT,
    checkpointer=MemorySaver(),   # ou SqliteSaver, PostgresSaver, ...
)

# Chaque conversation est identifiée par un thread_id :
config = {"configurable": {"thread_id": "user-42"}}
result = agent.invoke(
    {"messages": [HumanMessage(content="Salut")]},
    config=config,
)
\end{pythoncode}

Le checkpointer stocke l'état du graphe (donc l'historique) entre les
appels. L'utilisateur n'a plus à maintenir \texttt{history} côté client.

\begin{quote}\itshape
Pattern de production~: \texttt{MemorySaver} pour les tests, un
\texttt{Sqlite}- ou \texttt{Postgres}-saver pour la prod. Avec un
\texttt{thread\_id} par session utilisateur, on a un chat multi-utilisateur
sans aucune logique de mémoire côté client.
\end{quote}

\subsection{Vérification dans LangSmith}

Dans LangSmith, chaque \texttt{invoke} produit \textbf{une trace
distincte} (sauf si on utilise un \texttt{trace} explicite). Les traces
sont liées par leur projet et leur timestamp, mais pas automatiquement
groupées par conversation.

Pour grouper les runs d'un même utilisateur~:

\begin{itemize}[leftmargin=*]
  \item Avec checkpointer + \texttt{thread\_id}~: LangSmith expose le
        \texttt{thread\_id} dans les métadonnées de chaque run, on peut
        filtrer.
  \item Sans checkpointer~: passer manuellement un \texttt{run\_name} ou
        des \texttt{tags} pour grouper.
\end{itemize}

\begin{quote}\itshape
Astuce~: dans la version stateless (notre exercice), tu peux comparer
le \emph{nombre de tokens d'entrée} entre les runs successifs pour
visualiser la croissance de l'historique. Sur une conversation de 10
tours, c'est édifiant.
\end{quote}

\subsection*{Récapitulatif}
{\itshape\sloppy
\textit{Ce qu'il faut retenir :}
\begin{enumerate}[leftmargin=*]
  \item \texttt{create\_agent} est \textbf{stateless} par défaut~:
        chaque \texttt{.invoke()} est indépendant.
  \item Pour une mémoire conversationnelle simple~: accumuler les
        \texttt{messages} côté client, et réinjecter \emph{toute}
        l'historique à chaque tour.
  \item Le coût en tokens d'entrée \textbf{croît à chaque tour}. Sur
        conversations longues, tronquer ou résumer.
  \item Pour une persistence native côté serveur~: passer un
        \texttt{checkpointer} à \texttt{create\_agent} et identifier
        chaque session par un \texttt{thread\_id}.
  \item Pattern REPL standard~: init coûteux hors boucle, traitement
        rapide dans la boucle, terminaison propre sur Ctrl+C / EOF /
        ligne vide.
\end{enumerate}
\par}

\part{Pour aller plus loin}

% #####################################################################
% #####################################################################
\section{Ajouter un cinquième tool à l'agent}
% #####################################################################
% #####################################################################

\begin{tcolorbox}[title=Exercice bonus 1. Étendre l'inventaire de tools en une fonction]
Écris un tool \texttt{current\_python\_version()} qui renvoie la version
de Python utilisée par le processus, et ajoute-le à l'agent PyTutor.
Vérifie ensuite avec une question genre \emph{« Quelle version utilises-tu~?
Est-ce que je peux utiliser \texttt{match}/\texttt{case}~? »}.

\textbf{Contrainte}~: la docstring du tool doit mentionner les cas
d'usage typiques (compatibilité, features par version~: \texttt{match} en
3.10, \texttt{ExceptionGroup} en 3.11, etc.) pour que le LLM sache
\emph{quand} l'appeler.

\bigskip
\noindent\textbf{Objectifs pédagogiques}
\begin{itemize}[leftmargin=*]
  \item Mesurer la \textbf{trivialité} de l'ajout d'un tool une fois
        que l'agent est construit~: une fonction \texttt{@tool}, et
        on l'ajoute à la liste.
  \item Comprendre que le modèle découvre le nouveau tool
        \textbf{automatiquement} via son nom et sa docstring — pas de
        configuration explicite.
  \item Voir le rôle des cas d'usage dans la docstring~: c'est ce qui
        déclenche (ou pas) l'appel du tool selon la question.
\end{itemize}
\end{tcolorbox}

\subsection{Solution}

\subsubsection{Architecture}

\begin{center}
\begin{tikzpicture}[node distance=5mm and 10mm]
  \node[llm] (model) {Modèle LLM};

  \node[tool, below left=14mm and 22mm of model]  (t1) {search\_python\_docs};
  \node[tool, below left=14mm and  4mm of model]  (t2) {run\_python\_code};
  \node[tool, below right=14mm and  4mm of model] (t3) {generate\_exercise};
  \node[tool, below right=14mm and 22mm of model] (t4) {check\_pep8};
  \node[tool, below=24mm of model, fill=orange!15, draw=orange!70!black] (t5)
       {current\_python\_version\\\scriptsize \textbf{(nouveau)}};

  \draw[flow] (model) -- (t1);
  \draw[flow] (model) -- (t2);
  \draw[flow] (model) -- (t3);
  \draw[flow] (model) -- (t4);
  \draw[flow, draw=orange!70!black] (model) -- (t5) node[label, midway, right]{ajouté};
\end{tikzpicture}
\end{center}

Un cinquième tool s'ajoute sans modification du reste~: même
\texttt{system\_prompt}, même \texttt{model}, même architecture
LangGraph. Le modèle, à l'invocation, voit 5 tools dans son contrat au
lieu de 4.

\subsubsection{Le code complet}

\begin{pythoncode}
import sys
from langchain.agents import create_agent
from langchain_core.tools import tool
from langchain_core.messages import HumanMessage
from config import get_model
from phase3_agent import TOOLS, SYSTEM_PROMPT, print_trace


@tool
def current_python_version() -> str:
    """Renvoie la version Python du processus qui exécute PyTutor.

    À utiliser quand l'élève pose des questions sur la compatibilité, les
    features récentes (match/case en 3.10, ExceptionGroup en 3.11, etc.),
    ou simplement quand il demande quelle version est utilisée.
    """
    return f"Python {sys.version.split()[0]} ({sys.platform})"


def exercise_1() -> None:
    # On rajoute simplement le nouveau tool à la liste existante.
    agent = create_agent(
        model=get_model(),
        tools=TOOLS + [current_python_version],
        system_prompt=SYSTEM_PROMPT,
    )

    result = agent.invoke({
        "messages": [HumanMessage(content=(
            "Quelle version de Python utilises-tu pour exécuter mon code ? "
            "Est-ce que je peux utiliser le pattern matching avec `match` / `case` ?"
        ))]
    })
    print_trace(result["messages"])
\end{pythoncode}

\subsubsection{Explications pas à pas}

\subsubsection*{Étape 1 — La fonction Python décorée}

Un tool n'a besoin que de trois choses~: un décorateur \texttt{@tool},
une signature typée, une docstring~:

\begin{pythoncode}
@tool
def current_python_version() -> str:
    """Renvoie la version Python..."""
    return f"Python {sys.version.split()[0]} ({sys.platform})"
\end{pythoncode}

\textbf{Sans argument}~: la signature \texttt{() -> str} est minimale.
Le LLM n'a rien à fournir en entrée, juste à décider d'appeler le tool.

\subsubsection*{Étape 2 — L'ajout à la liste \texttt{tools}}

\begin{pythoncode}
tools=TOOLS + [current_python_version]
\end{pythoncode}

C'est tout. Aucun re-enregistrement, aucune configuration. Le nouveau
tool est passé à \texttt{create\_agent} qui l'injecte dans le contrat
du LLM.

\begin{quote}\itshape
\textbf{Pattern de composition}~: importer la liste \texttt{TOOLS} de
\texttt{phase3\_agent.py} et la combiner avec de nouveaux tools évite
de dupliquer la définition de l'agent. C'est l'équivalent d'une
\emph{extension} en orienté-objet, mais en plus simple~: juste de la
concaténation de listes.
\end{quote}

\subsubsection*{Étape 3 — La docstring déclenche l'usage}

La docstring liste des cas d'usage spécifiques~:
\begin{itemize}[leftmargin=*]
  \item « questions sur la compatibilité » — déclenche l'appel quand
        l'utilisateur demande s'il peut utiliser \texttt{walrus} ou autre
        feature récente.
  \item « features récentes (match/case en 3.10\dots) » — donne au LLM
        des \emph{exemples concrets} de questions qui doivent y mener.
  \item « ou simplement quand il demande quelle version » — couvre le
        cas direct.
\end{itemize}

\subsubsection*{Étape 4 — Le comportement attendu de l'agent}

Sur la question test :
\emph{« Quelle version de Python utilises-tu~? Est-ce que je peux
utiliser \texttt{match}~/~\texttt{case}~? »}

L'agent devrait~:
\begin{enumerate}[leftmargin=*]
  \item Appeler \texttt{current\_python\_version()}.
  \item Observer la sortie (par exemple \texttt{"Python 3.11.4 (darwin)"}).
  \item Répondre en mentionnant que \texttt{match/case} (PEP 634) est
        disponible depuis Python 3.10, et donc utilisable.
\end{enumerate}

\begin{quote}\itshape
Si l'agent n'appelle PAS le tool, c'est que la docstring n'est pas
assez explicite. Reformule-la pour ajouter le mot-clé qui correspond à
la question (par exemple, ajouter « pattern matching, match, case »
parmi les cas d'usage).
\end{quote}

\begin{tcolorbox}[
  colback=white,
  colframe=black!50,
  coltitle=black,
  fonttitle=\footnotesize\bfseries,
  fontupper=\footnotesize,
  title=Sortie de l'exécution,
  breakable,
]
\begin{verbatim}
======================================================================
EXERCICE 1 -- Ajout d'un tool current_python_version()
======================================================================

[USER]:
  Quelle version de Python utilises-tu pour exécuter mon code ?
  Est-ce que je peux utiliser le pattern matching avec
  `match` / `case` ?

[AI]:
  Je vais vérifier la version de Python utilisée.
  -> tool_use: current_python_version({})

[TOOL current_python_version]:
     Python 3.12.7 (darwin)

[AI]:
  Excellente nouvelle !

  Je utilise **Python 3.12.7**, ce qui signifie que **oui, tu peux
  tout à fait utiliser le pattern matching avec `match` / `case`** !

  Le pattern matching a été introduit en **Python 3.10**, donc tu as
  accès à cette fonctionnalité. Voici un petit exemple pour te
  montrer que ça fonctionne :

  -> tool_use: run_python_code({'code':
       'def classifier_nombre(n):\n'
       '    match n:\n'
       '        case 0: return "zéro"\n'
       '        case 1: return "un"\n'
       '        case 2: return "deux"\n'
       '        case _: return "autre nombre"\n'
       'print(classifier_nombre(0))\n'
       'print(classifier_nombre(1))\n'
       'print(classifier_nombre(5))'})

[TOOL run_python_code]:
     STDOUT:
     zéro
     un
     autre nombre
     Exit code: 0

[AI]:
  Parfait ! Comme tu vois, le pattern matching fonctionne sans
  problème. Tu peux utiliser :
   - Les **patterns simples** (valeurs littérales)
   - Les **patterns structurels** (tuples, listes, dictionnaires)
   - Les **patterns avec gardes** (conditions supplémentaires
     avec `if`)
   - Le **wildcard** `_` pour capturer n'importe quoi
\end{verbatim}
\end{tcolorbox}

\subsection{Vérification dans LangSmith}

L'arbre d'exécution doit montrer~:

\begin{center}
\begin{tikzpicture}[every node/.style={font=\footnotesize\ttfamily, anchor=west}]
  \node                              (a) {agent (LLM, tour 1)};
  \node[below=1mm of a, xshift=4mm]  (b) {Output : tool\_calls = [\{current\_python\_version\}]};
  \node[below=1mm of b, xshift=-4mm] (c) {tools (ToolNode)};
  \node[below=1mm of c, xshift=4mm]  (d) {current\_python\_version() $\rightarrow$ "Python 3.11.4..."};
  \node[below=1mm of d, xshift=-4mm] (e) {agent (LLM, tour 2)};
  \node[below=1mm of e, xshift=4mm]  (f) {Output : "J'utilise Python 3.11. match/case est..."};
\end{tikzpicture}
\end{center}

Vérifie aussi l'\emph{Input} du premier appel LLM~: les \textbf{5 tools}
doivent être listés dans le schéma envoyé, dont \texttt{current\_python\_version}
avec sa docstring.

\subsection*{Récapitulatif}
{\itshape\sloppy
\textit{Ce qu'il faut retenir :}
\begin{enumerate}[leftmargin=*]
  \item Ajouter un tool est \textbf{trivial}~: une fonction
        \texttt{@tool}, et on l'append à la liste passée à
        \texttt{create\_agent}.
  \item Pattern de composition~: \texttt{tools=TOOLS + [new\_tool]} évite
        la duplication.
  \item La docstring est le \textbf{seul} levier pour orienter le modèle
        sur le \emph{quand}-utiliser. Soigner les cas d'usage explicites.
  \item Si le modèle n'appelle pas le tool attendu~: c'est presque
        toujours un problème de docstring, pas de logique.
\end{enumerate}
\par}

% #####################################################################
% #####################################################################
\section{LangSmith pour observer un agent}
% #####################################################################
% #####################################################################

\begin{tcolorbox}[title=Exercice bonus 2. Tracer la boucle ReAct dans l'UI LangSmith]
Configure LangSmith (variables d'environnement \texttt{LANGCHAIN\_TRACING\_V2},
\texttt{LANGCHAIN\_API\_KEY}, \texttt{LANGCHAIN\_PROJECT}) et exécute
l'agent sur une question \textbf{multi-tool} pour obtenir une trace
riche. Ouvre la trace dans l'UI et explore l'arbre d'exécution.

\textbf{Contrainte}~: l'exercice doit vérifier programmatiquement la
présence des variables d'env et afficher un mode d'emploi clair si la
config manque.

\bigskip
\noindent\textbf{Objectifs pédagogiques}
\begin{itemize}[leftmargin=*]
  \item Configurer LangSmith pour le suivi d'un agent (variables
        spécifiques \texttt{LANGCHAIN\_*}).
  \item Lire l'arbre d'exécution d'un agent dans l'UI~: nœuds
        \texttt{agent} / \texttt{tools}, prompts exacts, tokens
        consommés, latence.
  \item Comprendre que \emph{l'agent n'a pas de mémoire magique}~: c'est
        l'historique qui s'accumule et est repassé au modèle à chaque
        tour. LangSmith le rend visible.
\end{itemize}
\end{tcolorbox}

\subsection{Solution}

\subsubsection{Architecture}

\begin{center}
\begin{tikzpicture}[node distance=5mm and 14mm]
  \node[data] (env) {.env\\\scriptsize LANGCHAIN\_TRACING\_V2=true\\\scriptsize LANGCHAIN\_API\_KEY=...\\\scriptsize LANGCHAIN\_PROJECT=PyTutor2};
  \node[block, below=of env] (agent) {agent.invoke(\dots)};

  \node[block, right=24mm of agent] (api) {API LangSmith};
  \node[store, above=of api] (cloud) {smith.langchain.com\\\scriptsize traces persistées};

  \draw[flowdata] (env)   -- (agent);
  \draw[flow]     (agent) -- (api) node[label, midway, above]{tous les nœuds tracés};
  \draw[flow]     (api)   -- (cloud);
\end{tikzpicture}
\end{center}

Quand \texttt{LANGCHAIN\_TRACING\_V2=true} et qu'une clé API valide est
présente, \emph{chaque} \texttt{Runnable} et \emph{chaque} appel LLM
sont automatiquement envoyés à LangSmith. Aucun changement de code
nécessaire.

\subsubsection{Le code complet}

\begin{pythoncode}
import os
from langchain_core.messages import HumanMessage
from phase3_agent import build_agent


def exercise_2_langsmith() -> None:
    # --- Vérification de la configuration ---
    tracing = os.getenv("LANGCHAIN_TRACING_V2", "").lower() in {"true", "1"}
    api_key = os.getenv("LANGCHAIN_API_KEY")
    project = os.getenv("LANGCHAIN_PROJECT", "default")

    print(f"LANGCHAIN_TRACING_V2 : {'V' if tracing else 'X'}")
    print(f"LANGCHAIN_API_KEY    : {'V' if api_key else 'X'}")
    print(f"LANGCHAIN_PROJECT    : {project}")

    if not (tracing and api_key):
        print("-> Configure ces variables dans ton .env puis relance.")
        return

    # --- Exécution d'un agent multi-tool pour une trace riche ---
    agent = build_agent()
    agent.invoke({
        "messages": [HumanMessage(content=(
            "Cherche dans la doc ce qu'est un dictionnaire en Python, "
            "puis génère un exercice débutant sur ce sujet, "
            "et vérifie que le code de référence est PEP 8 compliant."
        ))]
    })
    print(f"-> Trace envoyée. Ouvre https://smith.langchain.com (projet '{project}').")
\end{pythoncode}

\subsubsection{Explications pas à pas}

\subsubsection*{Étape 1 — Les variables d'environnement LangSmith}

\begin{center}
\begin{tabular}{p{5cm} p{9cm}}
\toprule
\textbf{Variable} & \textbf{Rôle} \\
\midrule
\texttt{LANGCHAIN\_TRACING\_V2}
& \texttt{"true"} pour activer l'envoi automatique des traces. Sans ça,
les chaînes/agents s'exécutent normalement, rien n'est envoyé. \\[1mm]
\texttt{LANGCHAIN\_API\_KEY}
& Clé personnelle créée sur \texttt{smith.langchain.com} (commence par
\texttt{lsv2\_pt\_\dots}). \\[1mm]
\texttt{LANGCHAIN\_PROJECT}
& Nom du projet où grouper les traces. Ex.~: \texttt{"PyTutor2"}. \\[1mm]
\texttt{LANGCHAIN\_ENDPOINT} & Optionnel.
\texttt{https://eu.api.smith.langchain.com} pour un compte EU. \\
\bottomrule
\end{tabular}
\end{center}

\begin{quote}\itshape
\textbf{Note de nommage}~: les variables historiques sont préfixées
\texttt{LANGCHAIN\_} ; la nouvelle version unifiée est préfixée
\texttt{LANGSMITH\_}. Les deux sont reconnues. Cohérence avec le reste
du projet recommandée.
\end{quote}

\subsubsection*{Étape 2 — Une question \textbf{multi-tool} pour une trace riche}

Pour avoir une trace pédagogiquement intéressante, choisis une question
qui force l'agent à enchaîner plusieurs tools~:

\begin{quote}\itshape
\emph{« Cherche dans la doc ce qu'est un dictionnaire en Python, puis
génère un exercice débutant sur ce sujet, et vérifie que le code de
référence est PEP 8 compliant. »}
\end{quote}

Cette question demande explicitement~:
\begin{enumerate}[leftmargin=*]
  \item \texttt{search\_python\_docs} (« cherche dans la doc »).
  \item \texttt{generate\_exercise} (« génère un exercice »).
  \item \texttt{check\_pep8} (« vérifie PEP 8 »).
\end{enumerate}

L'agent va donc faire \textbf{3 itérations} de la boucle ReAct, ce qui
donne une trace très lisible.

\subsubsection*{Étape 3 — Ce qu'on voit dans l'UI LangSmith}

Dans \texttt{smith.langchain.com}, dans ton projet~:

\begin{enumerate}[leftmargin=*]
  \item Une trace nommée \texttt{LangGraph} apparaît~: l'agent
        \emph{est} un graphe LangGraph.
  \item L'\textbf{arbre d'exécution} montre l'alternance
        \texttt{agent} (LLM) / \texttt{tools} (ToolNode) sur 3 cycles.
  \item Chaque \texttt{agent} (model call) expose~:
        \begin{itemize}
        \item \emph{Input} : système + historique + schémas des tools.
        \item \emph{Output} : la décision (tool\_calls) ou la réponse
              finale.
        \end{itemize}
  \item Chaque \texttt{tools} expose les noms des tools appelés, leurs
        arguments, et leurs sorties intégrales.
  \item Tokens et latence \emph{par appel} sont visibles, agrégés sur
        la trace racine.
\end{enumerate}

\subsubsection*{Étape 4 — L'insight clé : l'historique grossit}

\begin{quote}\itshape
À chaque tour du modèle, regarde l'onglet \emph{Input} dans la trace.
Tu vois \textbf{tout l'historique précédent} qui est repassé au LLM.
L'agent n'a pas de mémoire magique~: c'est juste l'historique qui
grossit, et que LangChain repasse à chaque appel. C'est aussi pourquoi
le coût en tokens d'entrée augmente à chaque tour.
\end{quote}

\subsection{Vérification dans LangSmith}

\emph{(Cet exercice EST la vérification.)}

Quatre choses à pointer dans la trace pour bien comprendre~:

\begin{itemize}[leftmargin=*]
  \item L'\textbf{alternance} \texttt{agent} $\to$ \texttt{tools} $\to$
        \texttt{agent} $\to$ \dots
  \item Le \emph{même} \texttt{system\_prompt} en tête de chaque appel
        \texttt{agent} (constant).
  \item L'\emph{historique qui grandit} dans l'input de chaque appel
        \texttt{agent}.
  \item La \textbf{dernière} étape : un \texttt{agent} qui produit un
        \texttt{Output} \textbf{sans} \texttt{tool\_calls} — c'est la
        réponse finale.
\end{itemize}

\subsection*{Récapitulatif}
{\itshape\sloppy
\textit{Ce qu'il faut retenir :}
\begin{enumerate}[leftmargin=*]
  \item 3 variables d'env~: \texttt{LANGCHAIN\_TRACING\_V2=true},
        \texttt{LANGCHAIN\_API\_KEY=\dots},
        \texttt{LANGCHAIN\_PROJECT=\dots}.
  \item Tracing automatique~: aucune modification de code une fois
        activé.
  \item Une trace d'agent montre l'\textbf{alternance}
        agent/tools/agent\dots\ jusqu'à la décision finale.
  \item L'historique grossit visible dans l'\emph{Input} de chaque
        appel LLM successif. C'est l'observation pédagogique
        fondamentale.
  \item LangSmith expose tokens et latence par appel, agrégés sur la
        trace racine~: indispensable pour optimiser un agent.
\end{enumerate}
\par}

% #####################################################################
% #####################################################################
\section{Tool avec paramètre \texttt{timeout} configurable}
% #####################################################################
% #####################################################################

\begin{tcolorbox}[title=Exercice bonus 3. Exposer un paramètre par défaut au modèle]
Modifie \texttt{run\_python\_code} pour accepter un paramètre
\texttt{timeout\_seconds: int = 5}. La docstring doit guider le modèle
sur \emph{quand} surcharger la valeur par défaut (calcul long manifeste)
et \emph{quand pas} (la majorité des cas). Teste avec deux scénarios~:
un code rapide et un code volontairement long.

\textbf{Contrainte}~: le timeout doit être \emph{clampé} côté serveur
($1 \leq t \leq 60$ secondes) pour empêcher des valeurs absurdes.

\bigskip
\noindent\textbf{Objectifs pédagogiques}
\begin{itemize}[leftmargin=*]
  \item Comprendre que les paramètres avec valeur par défaut sont
        \textbf{exposés au modèle} dans le schéma JSON du tool.
  \item Voir comment la docstring conditionne la décision du modèle
        d'override ou non.
  \item Mettre en place une \textbf{défense en profondeur}~: une
        consigne dans la docstring \emph{plus} un clamping côté code.
\end{itemize}
\end{tcolorbox}

\subsection{Solution}

\subsubsection{Architecture}

\begin{center}
\begin{tikzpicture}[node distance=5mm and 12mm]
  \node[llm] (model) {Modèle LLM};
  \node[data, below=of model] (schema)
       {Schéma du tool :\\\scriptsize \{code: str, timeout\_seconds: int = 5\}};
  \node[tool, below=of schema, fill=orange!8, draw=orange!70!black]
       (tool) {run\_python\_code\_v2\\\scriptsize clamp + subprocess};

  \node[block, below right=8mm and 14mm of tool] (cas1)
       {Cas 1 : code rapide\\\scriptsize $\rightarrow$ default (5s)};
  \node[block, below left=8mm and 14mm of tool] (cas2)
       {Cas 2 : calcul long\\\scriptsize $\rightarrow$ override (ex. 30s)};

  \draw[flowdata] (model) -- (schema);
  \draw[flow]     (schema)-- (tool);
  \draw[flow]     (tool)  -- (cas1);
  \draw[flow]     (tool)  -- (cas2);
\end{tikzpicture}
\end{center}

Le modèle voit la signature complète, lit la docstring, et décide à
chaque appel s'il garde la valeur par défaut ou la surcharge.

\subsubsection{Le code complet}

\begin{pythoncode}
import subprocess
import sys
import tempfile
from pathlib import Path

from langchain_core.tools import tool


@tool
def run_python_code_v2(code: str, timeout_seconds: int = 5) -> str:
    """Exécute du code Python et renvoie stdout + stderr.

    Args:
        code: Le code Python à exécuter, complet et autonome.
        timeout_seconds: Délai max avant interruption. Défaut : 5 secondes.
            Augmente cette valeur SEULEMENT si le code requiert manifestement
            un calcul long (entraînement, traitement de gros datasets, etc.).
            Maximum raisonnable : 60 secondes.
    """
    timeout_seconds = max(1, min(timeout_seconds, 60))  # clamp défensif

    with tempfile.NamedTemporaryFile(mode="w", suffix=".py", delete=False) as f:
        f.write(code)
        path = f.name
    try:
        result = subprocess.run(
            [sys.executable, path],
            capture_output=True, text=True, timeout=timeout_seconds,
        )
        parts = []
        if result.stdout:
            parts.append(f"STDOUT:\n{result.stdout}")
        if result.stderr:
            parts.append(f"STDERR:\n{result.stderr}")
        parts.append(f"Exit code: {result.returncode} (timeout utilisé: {timeout_seconds}s)")
        return "\n\n".join(parts)
    except subprocess.TimeoutExpired:
        return (
            f"ERREUR: timeout après {timeout_seconds} secondes. "
            "Si tu pensais que le code devait être long, augmente timeout_seconds. "
            "Sinon, il y a sans doute une boucle infinie ou une récursion non terminale."
        )
    finally:
        Path(path).unlink(missing_ok=True)
\end{pythoncode}

\subsubsection{Explications pas à pas}

\subsubsection*{Étape 1 — Le paramètre avec valeur par défaut}

\begin{pythoncode}
def run_python_code_v2(code: str, timeout_seconds: int = 5) -> str:
\end{pythoncode}

Le décorateur \texttt{@tool} extrait \emph{toute} la signature et la
traduit en JSON Schema~:

\begin{pythoncode}
{
    "type": "object",
    "properties": {
        "code": {"type": "string"},
        "timeout_seconds": {"type": "integer", "default": 5},
    },
    "required": ["code"],     # timeout_seconds N'EST PAS requis
}
\end{pythoncode}

Ce JSON Schema est envoyé au LLM. Le modèle peut donc \textbf{choisir}
de remplir \texttt{timeout\_seconds}, ou de laisser le tool utiliser le
défaut.

\subsubsection*{Étape 2 — La docstring guide la décision}

C'est la docstring du paramètre qui pilote le comportement~:

\begin{quote}\itshape
« Augmente cette valeur \textbf{SEULEMENT} si le code requiert
manifestement un calcul long (entraînement, traitement de gros datasets,
etc.). Maximum raisonnable~: 60 secondes. »
\end{quote}

Trois leviers~:
\begin{itemize}[leftmargin=*]
  \item \textbf{« SEULEMENT si »} en capitales~: signal fort que c'est
        une exception, pas la règle.
  \item \textbf{Exemples concrets}~: « entraînement, gros datasets ».
        Donne au LLM des templates de cas qui justifient l'override.
  \item \textbf{Maximum nommé}~: « 60 secondes »~: le LLM ne devrait
        pas dépasser cette borne, mais on a aussi un clamp côté serveur
        au cas où.
\end{itemize}

\subsubsection*{Étape 3 — Le clamp défensif côté Python}

\begin{pythoncode}
timeout_seconds = max(1, min(timeout_seconds, 60))
\end{pythoncode}

Pourquoi double protection (prompt + clamp)~?
\begin{itemize}[leftmargin=*]
  \item Le modèle peut \textbf{halluciner} une valeur (1000 secondes,
        ou même négative).
  \item Le modèle peut être \textbf{exploité}~: une question pourrait
        chercher à passer \texttt{timeout\_seconds=99999} pour bloquer
        le serveur.
  \item En production, on ne fait jamais confiance au LLM pour des
        valeurs critiques~: \textbf{défense en profondeur}.
\end{itemize}

\begin{quote}\itshape
La règle générale~: tout paramètre numérique exposé à un LLM doit avoir
des bornes \textbf{explicites} côté serveur. La docstring guide ; le
clamp protège.
\end{quote}

\subsubsection*{Étape 4 — Les deux cas attendus}

\textbf{Cas 1 — code rapide} :

\begin{pythoncode}
"Exécute ce code : print(sum(range(100)))"
\end{pythoncode}

L'agent appelle \texttt{run\_python\_code\_v2(code="\dots")} \textbf{sans}
spécifier \texttt{timeout\_seconds} — il utilise le défaut de 5s.
Dans l'historique, le \texttt{tool\_call} contient un seul argument~:
\texttt{\{"code": "\dots"\}}.

\textbf{Cas 2 — calcul long} (somme des premiers jusqu'à 5M) :

L'utilisateur précise « ça met une bonne dizaine de secondes ». L'agent
\emph{lit} cette information, comprend qu'il dépasse le défaut, et
appelle~:

\begin{pythoncode}
run_python_code_v2(code="...", timeout_seconds=30)
\end{pythoncode}

\begin{quote}\itshape
C'est \textbf{une preuve} que la docstring conditionne le comportement
réel du modèle. La phrase « Augmente cette valeur SEULEMENT si\dots » est
lue et appliquée. Tu peux tester en retirant cette phrase et observer
que le modèle, soit override systématiquement (overshoot), soit ne
l'override jamais (timeout au cas 2).
\end{quote}

\begin{tcolorbox}[
  colback=white,
  colframe=black!50,
  coltitle=black,
  fonttitle=\footnotesize\bfseries,
  fontupper=\footnotesize,
  title=Sortie de l'exécution,
  breakable,
]
\begin{verbatim}
======================================================================
EXERCICE 3 -- Tool avec paramètre timeout_seconds configurable
======================================================================

--- Cas 1 : code rapide (timeout par défaut attendu) ---

[USER]:
  Exécute ce code pour me dire la sortie :
  ```python
  print(sum(range(100)))
  ```

[AI]:
  -> tool_use: run_python_code_v2({'code':
                                    'print(sum(range(100)))'})
       (pas de timeout_seconds -> défaut 5s)

[TOOL run_python_code_v2]:
     STDOUT:
     4950
     Exit code: 0 (timeout utilisé: 5s)

[AI]:
  La sortie est **4950**.

  Ce code calcule la somme de tous les nombres de 0 à 99 (car
  `range(100)` génère les entiers de 0 à 99 inclus).
  Mathématiquement, c'est : 0 + 1 + 2 + ... + 99 = **4950**.


--- Cas 2 : calcul long, l'agent doit augmenter le timeout ---

[USER]:
  J'ai écrit un code qui calcule la somme des nombres premiers
  jusqu'à 5 millions. Ca met une bonne dizaine de secondes à tourner
  sur ma machine. Tu peux l'exécuter pour vérifier le résultat ?

  ```python
  def is_prime(n):
      if n < 2: return False
      for i in range(2, int(n**0.5)+1):
          if n % i == 0: return False
      return True

  total = sum(n for n in range(2, 5_000_000) if is_prime(n))
  print(total)
  ```

[AI]:
  Je vais exécuter ton code. Comme tu l'as mentionné, ça va prendre
  un peu de temps, donc je vais augmenter le délai d'exécution.
  -> tool_use: run_python_code_v2({
       'code': '...',
       'timeout_seconds': 30   <-- augmenté explicitement
     })

[TOOL run_python_code_v2]:
     STDOUT:
     838596693108
     Exit code: 0 (timeout utilisé: 30s)

[AI]:
  Parfait ! Ton code fonctionne correctement et donne le résultat :
  **838 596 693 108**

  Ton algorithme est correct, mais je peux te montrer comment
  l'optimiser significativement avec le **Crible d'Eratosthène**,
  qui trouve tous les nombres premiers jusqu'à une limite en une
  seule passe (probablement 10-20x plus rapide).

----------------------------------------------------------------------
OBSERVATIONS :
 - Au cas 1, l'agent appelle le tool SANS spécifier timeout_seconds
   (il garde la valeur par défaut 5s).
 - Au cas 2, l'agent voit dans ta question que le code prend ~10s,
   et passe explicitement timeout_seconds=30 (ou similaire).
 - C'est UNE preuve que la docstring conditionne le comportement :
   la phrase 'Augmente cette valeur SEULEMENT si...' est lue et
   appliquée par le modèle.
\end{verbatim}
\end{tcolorbox}

\subsection{Vérification dans LangSmith}

Dans la trace, examine l'onglet \emph{Tool calls} du \texttt{AIMessage}
qui invoque \texttt{run\_python\_code\_v2}~:

\begin{itemize}[leftmargin=*]
  \item \textbf{Cas 1}~: \texttt{args = \{"code": "\dots"\}}. La clé
        \texttt{timeout\_seconds} \textbf{n'apparaît pas} → défaut utilisé.
  \item \textbf{Cas 2}~: \texttt{args = \{"code": "\dots", "timeout\_seconds": 30\}}.
        L'override est explicite.
\end{itemize}

Le \texttt{ToolMessage} en sortie contient la trace exacte du timeout
utilisé (\texttt{"Exit code: 0 (timeout utilisé: 30s)"}), ce qui permet
de vérifier que le clamp côté serveur a bien laissé passer la valeur.

\subsection*{Récapitulatif}
{\itshape\sloppy
\textit{Ce qu'il faut retenir :}
\begin{enumerate}[leftmargin=*]
  \item Les \textbf{paramètres avec défaut} sont exposés au LLM dans le
        JSON Schema. Le LLM choisit de les remplir ou pas.
  \item La \texttt{description} du paramètre dans la docstring est ce
        qui guide le modèle pour \emph{quand} override.
  \item \textbf{Défense en profondeur}~: prompt qui guide + clamp côté
        serveur. Ne jamais faire confiance aux valeurs numériques
        produites par le LLM.
  \item Toujours répercuter la valeur effectivement utilisée dans la
        sortie du tool (« timeout utilisé : 30s »)~: ça aide l'agent
        à interpréter le résultat.
\end{enumerate}
\par}

% #####################################################################
% #####################################################################
\section{Tool qui appelle un sous-LLM (\texttt{evaluate\_student\_answer})}
% #####################################################################
% #####################################################################

\begin{tcolorbox}[title=Exercice bonus 4. Composer un agent et un sous-LLM dans un tool]
Crée un tool \texttt{evaluate\_student\_answer(question, answer)} qui
prend une question et la réponse d'un élève, et renvoie une évaluation
structurée (correct/incorrect, note, feedback, concepts manquants). Le
tool doit utiliser \textbf{en interne} une mini-chaîne LCEL avec
\texttt{with\_structured\_output(AnswerEvaluation)}.

\textbf{Contrainte}~: la mini-chaîne LCEL est construite au niveau du
module (une fois pour toute) pour éviter de la recréer à chaque appel
du tool.

\bigskip
\noindent\textbf{Objectifs pédagogiques}
\begin{itemize}[leftmargin=*]
  \item Maîtriser le pattern « tool qui appelle un autre LLM en
        interne »~: c'est la composition d'agents au sens large.
  \item Utiliser \texttt{with\_structured\_output} avec un schéma
        Pydantic riche pour produire des évaluations \emph{exploitables}.
  \item Comprendre l'architecture en \textbf{couches}~: agent principal
        $\to$ tool $\to$ sous-chaîne LCEL $\to$ sous-LLM.
  \item Observer le coût~: 2 appels LLM dans un seul tour d'agent.
\end{itemize}
\end{tcolorbox}

\subsection{Solution}

\subsubsection{Architecture}

\begin{center}
\begin{tikzpicture}[node distance=4mm and 8mm]
  \node[llm] (main) {Agent principal\\\scriptsize Claude (modèle 1)};
  \node[block, below=of main] (tcall) {tool\_call evaluate\_student\_answer(\dots)};
  \node[tool, below=of tcall] (tool) {evaluate\_student\_answer};
  \node[block, below=of tool] (chain) {chaîne LCEL\\\scriptsize \_evaluator\_chain};
  \node[llm, below=of chain] (sub) {Sous-LLM\\\scriptsize Claude (modèle 2 = même)};
  \node[schema, below=of sub] (sch) {AnswerEvaluation\\\scriptsize (Pydantic)};
  \node[data, below=of sch] (out) {string formaté\\\scriptsize « Correct : oui / Note : 7/10 / \dots »};

  \draw[flow] (main)  -- (tcall);
  \draw[flow] (tcall) -- (tool);
  \draw[flow] (tool)  -- (chain) node[label, midway, right]{\_evaluator\_chain.invoke(\dots)};
  \draw[flow] (chain) -- (sub);
  \draw[flow] (sub)   -- (sch);
  \draw[flow] (sch)   -- (out);

  \draw[flow, dashed, draw=orange!70!black]
       (out.east) ..controls +(5,0) and +(5,0).. (main.east)
       node[label, midway, right]{remonte à l'agent};
\end{tikzpicture}
\end{center}

Vu de l'agent principal, c'est un tool comme un autre. Vu de
l'intérieur, c'est un \textbf{deuxième appel LLM} caché.

\subsubsection{Le code complet}

\begin{pythoncode}
from pydantic import BaseModel, Field
from langchain_core.tools import tool
from langchain_core.prompts import ChatPromptTemplate
from config import get_model


class AnswerEvaluation(BaseModel):
    """Évaluation pédagogique d'une réponse d'élève."""

    is_correct: bool = Field(
        description="True si la réponse est essentiellement correcte, "
                    "False si elle contient une erreur substantielle."
    )
    score: int = Field(
        description="Note de 0 à 10. 0 = totalement faux, 10 = parfait. "
                    "Tiens compte de la précision technique ET de la clarté."
    )
    feedback: str = Field(
        description="Feedback pédagogique en 2-3 phrases pour l'élève. "
                    "Souligne ce qui est juste, puis ce qu'il faut corriger ou préciser."
    )
    missing_concepts: list[str] = Field(
        description="Concepts importants que la réponse aurait dû mentionner mais omet. "
                    "Liste vide si la réponse est complète."
    )


# La mini-chaîne LCEL utilisée à l'intérieur du tool. On la construit
# une seule fois en module-level pour éviter de la recréer à chaque appel.
_eval_prompt = ChatPromptTemplate.from_messages([
    ("system",
     "Tu es un évaluateur pédagogique de réponses Python. Sois exigeant "
     "techniquement mais bienveillant dans la formulation."),
    ("human",
     "Question posée à l'élève :\n{question}\n\n"
     "Réponse de l'élève :\n{answer}\n\n"
     "Évalue cette réponse."),
])

_evaluator_chain = (
    _eval_prompt
    | get_model().with_structured_output(AnswerEvaluation)
)


@tool
def evaluate_student_answer(question: str, answer: str) -> str:
    """Évalue la réponse d'un élève à une question technique sur Python.

    À utiliser quand un élève a écrit une réponse et que tu veux la noter
    objectivement plutôt que de te contenter d'un "oui" ou "non" qualitatif.
    Renvoie un score, un feedback, et la liste des concepts manquants.

    Args:
        question: La question posée à l'élève.
        answer: La réponse complète de l'élève à évaluer.
    """
    # Sous-LLM : on appelle une mini-chaîne LCEL avec sortie structurée.
    result = _evaluator_chain.invoke({"question": question, "answer": answer})

    # On reformate l'objet Pydantic en string pour l'agent.
    # (Les tools renvoient toujours du texte, c'est leur protocole.)
    missing = ", ".join(result.missing_concepts) if result.missing_concepts else "(aucun)"
    return (
        f"Correct : {'oui' if result.is_correct else 'non'}\n"
        f"Note : {result.score}/10\n"
        f"Concepts manquants : {missing}\n"
        f"Feedback : {result.feedback}"
    )
\end{pythoncode}

\subsubsection{Explications pas à pas}

\subsubsection*{Étape 1 — Le schéma Pydantic d'évaluation}

\texttt{AnswerEvaluation} contient 4 champs orientés
\emph{pédagogie}~: un booléen pour le verdict net, un score chiffré, un
feedback en langage naturel, et une liste de concepts manquants.

Le choix des types n'est pas neutre~:
\begin{itemize}[leftmargin=*]
  \item \texttt{is\_correct: bool}~: verdict binaire utile en aval
        (par exemple pour filtrer les bonnes réponses).
  \item \texttt{score: int}~: granularité 0-10, exploitable pour
        agrégation et stats.
  \item \texttt{feedback: str}~: texte à montrer à l'élève.
  \item \texttt{missing\_concepts: list[str]}~: liste itérable, pas
        une chaîne JSON. Permet de calculer du \emph{coverage} sur
        plusieurs réponses.
\end{itemize}

\subsubsection*{Étape 2 — La mini-chaîne LCEL construite \emph{une fois}}

\begin{pythoncode}
_eval_prompt = ChatPromptTemplate.from_messages([...])
_evaluator_chain = _eval_prompt | get_model().with_structured_output(AnswerEvaluation)
\end{pythoncode}

\textbf{Pourquoi au niveau du module~?} \texttt{get\_model()} crée une
instance du client (charge la config, ouvre éventuellement des
connexions). On la construit une fois à l'import et on la réutilise.
Sinon, chaque appel du tool recréerait l'instance — coûteux et
inutile.

\begin{quote}\itshape
Pattern important~: tout ce qui est \emph{coûteux à construire et
réutilisable} (modèle, chaîne, retriever, embeddings\dots) vit au
niveau du module. Les fonctions et tools ne font que \texttt{.invoke()}
sur ces objets.
\end{quote}

\subsubsection*{Étape 3 — Le tool reformate en string}

Le tool reçoit deux strings (question, answer), invoque la chaîne, et
\textbf{reformate} l'objet Pydantic en string~:

\begin{pythoncode}
result = _evaluator_chain.invoke({"question": question, "answer": answer})
# result est un objet AnswerEvaluation

return (
    f"Correct : {'oui' if result.is_correct else 'non'}\n"
    f"Note : {result.score}/10\n"
    ...
)
\end{pythoncode}

\textbf{Pourquoi reformater~?} Le protocole des tools est~: \emph{ils
renvoient une string}. Cette string devient le contenu d'un
\texttt{ToolMessage} dans l'historique, ensuite \emph{relu} par l'agent
principal pour formuler sa réponse à l'utilisateur.

\subsubsection*{Étape 4 — L'architecture en couches}

Quand l'agent principal traite la question \emph{« Évalue cette réponse
d'élève »}, l'arbre d'exécution a quatre niveaux~:

\begin{enumerate}[leftmargin=*]
  \item \textbf{Agent principal (Claude)}~: décide d'appeler
        \texttt{evaluate\_student\_answer(question, answer)}.
  \item \textbf{Tool \texttt{evaluate\_student\_answer}}~: Python exécute
        la fonction.
  \item \textbf{Sous-LLM (Claude à nouveau)}~: appelé en interne par
        la mini-chaîne. Reçoit son \emph{propre} prompt système et
        renvoie un \texttt{AnswerEvaluation} validé par Pydantic.
  \item \textbf{Le tool reformate} l'objet en string et la rend à
        l'agent principal.
  \item \textbf{L'agent principal} intègre le résultat et formule sa
        réponse à l'utilisateur.
\end{enumerate}

C'est \textbf{2 appels LLM} dans un seul tour d'agent~: un par
l'agent principal pour décider, un par le sous-LLM pour évaluer.

\subsubsection*{Étape 5 — Pourquoi ce pattern est important}

Le pattern « tool qui appelle un sous-LLM » résout des problèmes
fréquents~:

\begin{itemize}[leftmargin=*]
  \item \textbf{Compétence spécialisée}~: l'agent principal n'est pas un
        évaluateur ; il \emph{délègue} à un LLM dédié avec son prompt
        d'évaluation.
  \item \textbf{Sortie structurée garantie}~: l'agent principal manipule
        \texttt{AnswerEvaluation} via le tool, donc \texttt{score} est
        toujours un \texttt{int} entre 0 et 10, jamais une string
        ambiguë.
  \item \textbf{Réutilisabilité}~: la mini-chaîne peut être appelée
        directement ailleurs (interface enseignant) sans passer par
        l'agent.
\end{itemize}

\begin{tcolorbox}[
  colback=white,
  colframe=black!50,
  coltitle=black,
  fonttitle=\footnotesize\bfseries,
  fontupper=\footnotesize,
  title=Sortie de l'exécution,
  breakable,
]
\begin{verbatim}
======================================================================
EXERCICE 4 -- Tool 'evaluate_student_answer' (sous-LLM caché)
======================================================================

[USER]:
  Un élève a répondu à une question d'examen. Evalue-la pour moi et
  formule un commentaire que je pourrais lui envoyer.

  Question : A quoi sert le mot-clé `yield` en Python et en quoi
  diffère-t-il de `return` ?

  Réponse de l'élève : << `yield` permet de retourner plusieurs
  valeurs depuis une fonction, contrairement à `return` qui n'en
  retourne qu'une seule. >>

[AI]:
  Je vais évaluer cette réponse de manière objective.
  -> tool_use: evaluate_student_answer({
       'question': 'A quoi sert le mot-clé `yield`...',
       'answer':   '`yield` permet de retourner...'
     })

[TOOL evaluate_student_answer]:
     Correct : non
     Note : 4/10
     Concepts manquants :
       - Générateurs et itérateurs
       - Lazy evaluation (évaluation paresseuse)
       - Sauvegarde de l'état d'exécution avec yield
       - Efficacité mémoire des générateurs
       - Différence entre retourner une liste et utiliser yield
     Feedback : Tu as identifié une différence importa...

[AI]:
  ## Commentaire à envoyer à l'élève :

  **Evaluation : 4/10**

  Bonne intuition, mais ta réponse est **incomplète et imprécise**.
  Voici ce qui manque :

  ### Ce que tu as bien compris
  Tu as raison que `yield` permet de retourner plusieurs valeurs, ce
  qui est une différence avec `return`.

  ### Ce qui manque
  La vraie différence entre `yield` et `return` ne porte pas
  seulement sur le **nombre** de valeurs, mais sur **comment** et
  **quand** elles sont retournées :

  1. **`yield` crée un générateur** : une fonction qui retourne les
     valeurs **une à une, à la demande** (lazy evaluation), sans
     tout calculer d'un coup.
  2. **`return` retourne tout immédiatement** : la fonction s'arrête
     et retourne une valeur unique (ou une liste si tu veux
     plusieurs valeurs).
  3. **Avantage mémoire** : avec `yield`, tu économises de la
     mémoire car les valeurs ne sont générées que quand on les
     demande.

  ### Exemple pour clarifier
  ```python
  # Avec return : tout est calculé et stocké en mémoire
  def nombres_return():
      return [1, 2, 3, 4, 5]

  # Avec yield : les valeurs sont générées à la demande
  def nombres_yield():
      yield 1; yield 2; yield 3; yield 4; yield 5
  ```

  **A revoir** : les concepts de **générateurs**, **itérateurs** et
  **lazy evaluation** pour bien maîtriser `yield`.

----------------------------------------------------------------------
ARCHITECTURE EN COUCHES A VISUALISER :
 1. AGENT PRINCIPAL (Claude) : décide d'appeler evaluate_student_answer.
 2. TOOL evaluate_student_answer : Python exécute la fonction.
 3. SOUS-LLM (Claude à nouveau) : appelé en interne par la mini-chaîne.
    Reçoit son propre prompt, renvoie un objet Pydantic structuré.
 4. Le tool reformate l'objet en string et le rend à l'agent.
 5. L'agent intègre le résultat et formule sa réponse à l'utilisateur.

 -> 2 appels LLM dans ce tour, pour 1 invocation d'agent.
 -> Si tu actives LangSmith, tu verras les 2 traces emboîtées.
\end{verbatim}
\end{tcolorbox}

\subsection{Vérification dans LangSmith}

L'arbre d'exécution montre les \textbf{traces emboîtées}~:

\begin{center}
\begin{tikzpicture}[every node/.style={font=\footnotesize\ttfamily, anchor=west}]
  \node                              (a) {agent (LLM \#1) ← Claude décide};
  \node[below=1mm of a]              (b) {  Output: tool\_calls = [evaluate\_student\_answer]};
  \node[below=1mm of b]              (c) {tools (ToolNode)};
  \node[below=1mm of c, xshift=4mm]  (d) {evaluate\_student\_answer(question, answer)};
  \node[below=1mm of d, xshift=4mm]  (e) {RunnableSequence (\_evaluator\_chain)};
  \node[below=1mm of e, xshift=4mm]  (f) {ChatPromptTemplate};
  \node[below=1mm of f]              (g) {ChatAnthropic (LLM \#2) $\rightarrow$ AnswerEvaluation};
  \node[below=1mm of g, xshift=-12mm](h) {agent (LLM \#3) ← Claude répond};
\end{tikzpicture}
\end{center}

Tu peux cliquer sur chaque \texttt{ChatAnthropic} pour voir les
prompts \emph{différents}~: l'un est le system\_prompt de l'agent
principal, l'autre est celui de l'évaluateur.

\subsection*{Récapitulatif}
{\itshape\sloppy
\textit{Ce qu'il faut retenir :}
\begin{enumerate}[leftmargin=*]
  \item Pattern \textbf{tool $\to$ sous-LLM}~: un tool peut contenir
        une chaîne LCEL complète, avec son propre prompt et sortie
        structurée.
  \item Construire les objets coûteux (modèle, chaîne) au niveau du
        \textbf{module}, pas dans la fonction du tool.
  \item Les tools renvoient toujours du \textbf{texte}~: reformater
        l'objet Pydantic en string compréhensible pour l'agent.
  \item Coût~: \textbf{2 appels LLM} dans un seul tour d'agent. Visible
        en clair dans LangSmith (traces emboîtées).
  \item Cas d'usage~: déléguer une compétence spécialisée
        (évaluation, classification, extraction) tout en gardant un
        agent généraliste en orchestrateur.
\end{enumerate}
\par}

% #####################################################################
% #####################################################################
\section{Streaming avec \texttt{agent.stream()}}
% #####################################################################
% #####################################################################

\begin{tcolorbox}[title=Exercice bonus 5. Streamer un agent en deux modes complémentaires]
Passe de \texttt{agent.invoke()} à \texttt{agent.stream()} pour obtenir
une UX réactive. Expérimente avec deux modes~:
\texttt{stream\_mode="updates"} (chaque étape du graphe émet un chunk au
moment où le nœud termine) et \texttt{stream\_mode="messages"} (chaque
token du modèle est émis individuellement, effet machine à écrire).

\textbf{Contrainte}~: dans le mode \texttt{messages}, filtrer pour ne
streamer que les tokens du nœud \texttt{agent} (pas ceux d'éventuels
sous-LLM dans des tools).

\bigskip
\noindent\textbf{Objectifs pédagogiques}
\begin{itemize}[leftmargin=*]
  \item Connaître les trois modes de \texttt{agent.stream()}~:
        \texttt{updates}, \texttt{messages}, \texttt{values}.
  \item Savoir lire un chunk \texttt{updates}~: c'est un dict
        \texttt{\{node\_name: state\_diff\}}.
  \item Savoir filtrer le streaming \texttt{messages} par nœud (via la
        \texttt{metadata}).
  \item Combiner les deux modes pour une UX moderne (badges d'étapes +
        machine à écrire).
\end{itemize}
\end{tcolorbox}

\subsection{Solution}

\subsubsection{Architecture}

\begin{center}
\resizebox{\linewidth}{!}{%
\begin{tikzpicture}[node distance=5mm and 10mm]
  \node[block] (agent) {agent.stream(\{messages: [\dots]\}, stream\_mode=\dots)};

  \node[block, below left=14mm and 14mm of agent, fill=teal!8] (updates)
       {stream\_mode="updates"\\\scriptsize chunks par nœud terminé};
  \node[block, below right=14mm and 14mm of agent, fill=teal!8] (messages)
       {stream\_mode="messages"\\\scriptsize chunks token par token};

  \node[data, below=of updates] (u_use)
       {Badges : « utilise X »\\« utilise Y »\\\dots};
  \node[data, below=of messages] (m_use)
       {Effet machine à écrire\\(comme ChatGPT)};

  \draw[flow] (agent) -- (updates);
  \draw[flow] (agent) -- (messages);
  \draw[flow] (updates) -- (u_use);
  \draw[flow] (messages) -- (m_use);
\end{tikzpicture}%
}
\end{center}

Les deux modes \textbf{ne sont pas exclusifs}~: en production, on les
combine. \texttt{updates} pour afficher des badges au-dessus du chat
(« l'agent utilise \texttt{search\_python\_docs}\dots ») ;
\texttt{messages} pour le texte qui s'écrit progressivement.

\subsubsection{Le code complet}

\begin{pythoncode}
from langchain_core.messages import HumanMessage
from phase3_agent import build_agent


def exercise_5_updates() -> None:
    """Streaming au niveau des ÉTAPES (mode 'updates')."""
    agent = build_agent()

    for chunk in agent.stream(
        {"messages": [HumanMessage(content=(
            "Cherche dans la doc ce qu'est un itérable, "
            "puis génère un exercice débutant sur ce sujet."
        ))]},
        stream_mode="updates",
    ):
        # chunk est un dict {nom_du_noeud: {clés_modifiées: valeurs}}.
        # Avec create_agent, on a deux noeuds principaux : 'agent' et 'tools'.
        for node_name, node_update in chunk.items():
            messages = node_update.get("messages", [])
            for msg in messages:
                if hasattr(msg, "tool_calls") and msg.tool_calls:
                    for tc in msg.tool_calls:
                        print(f"  [{node_name}] -> appel {tc['name']}({list(tc['args'].keys())})")
                elif hasattr(msg, "name") and msg.name:
                    # ToolMessage : a un .name pointant vers le tool qui a renvoyé
                    preview = (msg.content or "")[:80].replace("\n", " ")
                    print(f"  [{node_name}] <- retour de {msg.name}: {preview}...")
                elif msg.content:
                    preview = msg.content[:120].replace("\n", " ")
                    print(f"  [{node_name}] >> {preview}...")


def exercise_5_tokens() -> None:
    """Streaming token par token (mode 'messages')."""
    agent = build_agent()

    # Mode 'messages' : on reçoit des tuples (chunk_message, metadata).
    for chunk, metadata in agent.stream(
        {"messages": [HumanMessage(content=(
            "Explique-moi en deux phrases ce qu'est une compréhension de liste."
        ))]},
        stream_mode="messages",
    ):
        # On ne stream que le texte produit par le noeud 'agent' (pas les tools)
        if metadata.get("langgraph_node") == "agent" and chunk.content:
            text = chunk.content if isinstance(chunk.content, str) else "".join(
                b.get("text", "") for b in chunk.content if isinstance(b, dict)
            )
            print(text, end="", flush=True)
    print()
\end{pythoncode}

\subsubsection{Explications pas à pas}

\subsubsection*{Étape 1 — Les trois modes de streaming}

\begin{center}
\begin{tabular}{p{2.5cm} p{11.5cm}}
\toprule
\textbf{Mode} & \textbf{Ce qu'on reçoit} \\
\midrule
\texttt{updates}
& Un chunk par \textbf{mise à jour} de l'état du graphe (= un nœud
termine). Format~: \texttt{\{node\_name: state\_diff\}}. Idéal pour
afficher « l'agent a appelé tel tool » en temps réel. \\[1mm]
\texttt{messages}
& Un chunk par \textbf{token} produit par un modèle. Format~: tuple
\texttt{(chunk\_msg, metadata)}. Idéal pour l'effet machine à écrire. \\[1mm]
\texttt{values}
& L'\textbf{état complet} du graphe à chaque étape (très verbeux, utile
en debug). \\
\bottomrule
\end{tabular}
\end{center}

\subsubsection*{Étape 2 — Mode \texttt{updates}~: chunks par étape}

\begin{pythoncode}
for chunk in agent.stream({"messages": [...]}, stream_mode="updates"):
    for node_name, node_update in chunk.items():
        # node_name est typiquement 'agent' ou 'tools'
        messages = node_update.get("messages", [])
        # messages contient ce que ce nœud A AJOUTÉ à l'état
\end{pythoncode}

Avec \texttt{create\_agent}, deux nœuds principaux émettent des updates~:
\begin{itemize}[leftmargin=*]
  \item \texttt{agent}~: produit des \texttt{AIMessage} (avec ou sans
        \texttt{tool\_calls}).
  \item \texttt{tools}~: produit des \texttt{ToolMessage} (résultats des
        tools).
\end{itemize}

Dans la boucle, on inspecte le contenu pour afficher~: appel de tool,
retour de tool, ou réponse textuelle.

\subsubsection*{Étape 3 — Mode \texttt{messages}~: tokens du LLM}

\begin{pythoncode}
for chunk, metadata in agent.stream({"messages": [...]}, stream_mode="messages"):
    if metadata.get("langgraph_node") == "agent" and chunk.content:
        print(chunk.content, end="", flush=True)
\end{pythoncode}

Deux subtilités~:

\begin{itemize}[leftmargin=*]
  \item Le chunk est un \texttt{AIMessageChunk} (ou un autre type de
        message). \texttt{chunk.content} peut être une \texttt{str} ou
        une liste de blocs (selon la version de l'API). Le code gère
        les deux cas.
  \item Le \texttt{metadata.langgraph\_node} indique de \emph{quel}
        nœud du graphe provient le token. On filtre sur
        \texttt{"agent"} pour ne pas streamer accidentellement les
        tokens d'un éventuel sous-LLM caché dans un tool.
\end{itemize}

\subsubsection*{Étape 4 — \texttt{print(\dots, end="", flush=True)} : le secret}

Pour que les tokens s'écrivent en temps réel sur le terminal~:
\begin{itemize}[leftmargin=*]
  \item \texttt{end=""}~: pas de saut de ligne après chaque chunk.
  \item \texttt{flush=True}~: force l'écriture immédiate (sans ça,
        Python bufferise par ligne ou par bloc de 4ko).
\end{itemize}

C'est le même pattern que pour le streaming \emph{chains} de la Phase 1.

\subsubsection*{Étape 5 — Combiner les deux modes en production}

Dans une vraie UI~:
\begin{itemize}[leftmargin=*]
  \item En \textbf{haut} du chat~: badges « l'agent réfléchit\dots »,
        « utilise \texttt{search\_python\_docs}\dots »~— issus du mode
        \texttt{updates}.
  \item Dans la \textbf{bulle de réponse}~: texte qui s'écrit
        progressivement~— issu du mode \texttt{messages}.
\end{itemize}

LangChain permet d'obtenir les deux en parallèle avec
\texttt{stream\_mode=["updates", "messages"]}~:

\begin{pythoncode}
for kind, chunk in agent.stream(
    {"messages": [...]},
    stream_mode=["updates", "messages"],   # mode multiple
):
    if kind == "updates":
        # afficher badge
        pass
    elif kind == "messages":
        # afficher token
        pass
\end{pythoncode}

\begin{quote}\itshape
\textbf{Règle pratique}~: en production, toujours streamer. Une UX qui
attend 8 secondes en boîte noire avant d'afficher la réponse complète
est perçue comme cassée — même si la durée totale est identique au
streaming.
\end{quote}

\begin{tcolorbox}[
  colback=white,
  colframe=black!50,
  coltitle=black,
  fonttitle=\footnotesize\bfseries,
  fontupper=\footnotesize,
  title=Sortie de l'exécution,
  breakable,
]
\begin{verbatim}
======================================================================
EXERCICE 5a -- Streaming par étapes (stream_mode='updates')
======================================================================

Question envoyée. Affichage des étapes au fil de l'eau :

  [model] -> appel search_python_docs(['query'])
  [model] -> appel generate_exercise(['topic', 'difficulty'])
[OK] Rechargement de l'index depuis ./chroma_db
  [tools] <- retour de search_python_docs: Source:
            https://docs.python.org/3/glossary.html or the script
            being run has fini...
  [tools] <- retour de generate_exercise: Exercice sur 'itérables'
            (niveau débutant) : Code de référence : # Exemple 1 :...
  [model] [chat] ## Définition d'un itérable
            D'après la **documentation officielle Python**, un
            **itérable** est :
            > Un objet capable ...

======================================================================
EXERCICE 5b -- Streaming token par token (stream_mode='messages')
======================================================================

Réponse qui s'affiche au fur et à mesure (effet machine à écrire) :

  [... tokens consommés un par un, sortie typewriter ...]

----------------------------------------------------------------------
QUAND UTILISER QUEL MODE :
 - 'updates'  : UI qui montre 'l'agent réfléchit / l'agent utilise X'.
                Idéal pour les workflows long-running où
                l'utilisateur attendrait sinon dans le vide.
 - 'messages' : effet machine à écrire, comme ChatGPT et Claude.ai.
                Indispensable pour une UX de chat moderne.
 - On combine souvent les deux : on stream les tokens du texte
   final, et on affiche des badges 'utilise tool X' pour les
   étapes intermédiaires.
\end{verbatim}
\end{tcolorbox}

\subsection{Vérification dans LangSmith}

\texttt{agent.stream()} produit la \textbf{même trace} que
\texttt{agent.invoke()} dans LangSmith~: les nœuds sont tracés au moment
où ils terminent. La seule différence est côté client~: on consomme
incrémentalement au lieu d'attendre la fin.

\begin{quote}\itshape
LangSmith expose la métrique \emph{Time to first token (TTFT)} dans les
runs streamés. Sur un agent à 2 itérations ReAct, le TTFT est de
l'ordre du \texttt{search\_python\_docs} (rapide) ou du premier
\texttt{ChatAnthropic} (variable). Une UI streamée affiche le premier
contenu en $\sim$300 ms au lieu de $\sim$5 s.
\end{quote}

\subsection*{Récapitulatif}
{\itshape\sloppy
\textit{Ce qu'il faut retenir :}
\begin{enumerate}[leftmargin=*]
  \item \texttt{agent.stream(input, stream\_mode=\dots)} a 3 modes :
        \texttt{updates} (par nœud), \texttt{messages} (par token),
        \texttt{values} (état complet).
  \item Mode \texttt{updates}~: chunks \texttt{\{node\_name: state\_diff\}},
        idéal pour les badges « l'agent fait X ».
  \item Mode \texttt{messages}~: tuples \texttt{(chunk\_msg, metadata)}.
        Filtrer sur \texttt{metadata["langgraph\_node"] == "agent"}
        pour exclure les sous-LLM.
  \item \texttt{print(\dots, end="", flush=True)} reste la condition
        sine qua non du streaming visible côté terminal.
  \item En production, on \textbf{combine} les deux modes
        (\texttt{stream\_mode=["updates", "messages"]}) pour une UX
        moderne avec badges d'étapes et machine à écrire.
\end{enumerate}
\par}


\end{document}
